% main.tex — پایان‌نامه: درجه‌بندی خودکار رتینوپاتی دیابتی و ادم ماکولا
% ساخت:  cd thesis && xelatex main && bibtex main && xelatex main && xelatex main
% یا:    make -C thesis
%
% همهٔ عددهای فارسی متن از docs/generated/thesis_numbers.json می‌آیند و با
% «python3 tools/farsi_lint.py --numbers thesis/*.tex» بررسی می‌شوند.

\documentclass[12pt,a4paper,oneside]{report}

\usepackage{amsmath,amssymb}
\usepackage{graphicx}
\usepackage{float}
\usepackage{longtable}
\usepackage{array}
\usepackage{booktabs}
\usepackage[table]{xcolor}
\usepackage[margin=2.7cm,bindingoffset=0.6cm]{geometry}
\usepackage{setspace}
\usepackage{enumitem}
\usepackage{pifont}   % \ding{55} — نشانهٔ رد در جدول‌های بخش نخست
\usepackage[hidelinks]{hyperref}

% ---- فارسی. xepersian باید پس از بقیهٔ بسته‌ها بارگذاری شود.
\usepackage{xepersian}

% قلم فارسی: وزیرمتن، با وزن سیاه واقعی (نه سیاهِ ساختگی) و ایتالیک که برای فارسی
% تعریف نمی‌شود، پس \emph{} با همان وزن معمولی حروف‌چینی می‌شود.
\settextfont[Path=fonts/,
             BoldFont=Vazirmatn-Bold.ttf,
             ItalicFont=Vazirmatn-Medium.ttf,
             BoldItalicFont=Vazirmatn-Bold.ttf,
             Scale=1.0]{Vazirmatn-Regular.ttf}
% قلم لاتین: با نام پرونده بارگذاری می‌شود، نه با نام خانواده — نام خانواده تنها وقتی
% کار می‌کند که قلم در سامانه نصب باشد، و اینجا از خودِ توزیع تک می‌آید.
\setlatintextfont[Extension=.otf,
                  UprightFont=lmroman10-regular,
                  BoldFont=lmroman10-bold,
                  ItalicFont=lmroman10-italic,
                  BoldItalicFont=lmroman10-bolditalic,
                  Scale=0.92]{lmroman10-regular}
\setdigitfont[Path=fonts/,
              BoldFont=Vazirmatn-Bold.ttf,
              ItalicFont=Vazirmatn-Medium.ttf,
              BoldItalicFont=Vazirmatn-Bold.ttf]{Vazirmatn-Regular.ttf}

\linespread{1.35}
\setlength{\parskip}{0.45em}

% محیط فهرست برای نشانه‌ها و کوته‌نوشت‌ها
\newenvironment{latinlist}
  {\begin{list}{}{\setlength{\labelwidth}{4.2em}%
    \setlength{\leftmargin}{5.0em}\setlength{\itemsep}{0.25em}%
    \renewcommand{\makelabel}[1]{\textbf{##1}\hfill}}}
  {\end{list}}

% \wrongnum{...} برای عددی که متن عمداً به شکل نادرست نقل می‌کند (بخش شکست تبار داده).
% ویراستار خودکار همین نشانه را می‌بیند و آن عدد را از بررسی کنار می‌گذارد، پس استثنا
% در متن دیده می‌شود نه در قاعدهٔ آزمون.
\newcommand{\wrongnum}[1]{\LR{#1}}

% \num{...} هر عدد فارسی را در یک بازهٔ صریح چپ‌به‌راست می‌گذارد.
% دلیلش این است که پیاده‌سازی دوسویهٔ زی‌تک، پیاده‌سازی کامل الگوریتم یونیکد نیست:
% عددی مثل «۰٫۵» که به ترتیب منطقی نوشته شود، بدون این پوشش «۵٫۰» چاپ می‌شود، چون
% جداکنندهٔ اعشار دو دستهٔ رقم را به هم نمی‌بندد و دسته‌ها مثل هر محتوای راست‌به‌چپ
% دیگری از راست به چپ چیده می‌شوند. با \LR درست چاپ می‌شود؛ با \lr عدد از صفحه
% \emph{ناپدید} می‌شود، چون \lr قلم لاتین را می‌آورد و آن قلم رقم فارسی ندارد.
% هر دو حالت با چاپ آزمایشی بررسی شده‌اند، نه با حدس.
\newcommand{\num}[1]{\LR{#1}}

\hypersetup{pdftitle={درجه‌بندی خودکار رتینوپاتی دیابتی و ادم ماکولای دیابتی},
            pdflang={fa}}

\begin{document}

% ------------------------------------------------------------------ صفحهٔ عنوان
\begin{titlepage}
\centering
\vspace*{2cm}
{\Large به نام خدا}\\[2.5cm]
{\huge\bfseries درجه‌بندی خودکار رتینوپاتی دیابتی\\[0.4em]
و ادم ماکولای دیابتی از تصویر ته‌چشم}\\[1.2em]
{\Large قاعدهٔ تصمیم، سقف نظارت، و تبار داده}\\[3cm]
{\large پایان‌نامه}\\[2.5cm]
{\large علیرضا نوروزی}\\[1cm]
{\large شهریور \num{۱۴۰۵}}
\vfill
\end{titlepage}

\cleardoublepage
\pagenumbering{roman}

% پیش‌مطالب — چکیدهٔ فارسی، پیش‌گفتار، فهرست نشانه‌ها و کوته‌نوشت‌ها
% همهٔ عددها از docs/generated/thesis_numbers.json می‌آیند.

\chapter*{چکیده}
\addcontentsline{toc}{chapter}{چکیده}
\markboth{چکیده}{چکیده}

رتینوپاتی دیابتی بزرگ‌ترین علت نابینایی پیشگیری‌پذیر در جمعیت شاغل جهان است و ادم ماکولای
دیابتی شایع‌ترین سازوکار از دست رفتن بینایی مرکزی. درجه‌بندی خودکار تصویر ته‌چشم اعداد بالایی
گزارش می‌کند، در حالی که تا پایان آوریل \num{۲۰۲۴} تنها \textbf{\num{۳}} دستگاه هوش مصنوعی چشم‌پزشکی مجوز
سازمان غذا و داروی ایالات متحده را داشته‌اند. این پژوهش از همین شکاف آغاز می‌کند و می‌پرسد
عددهای گزارش‌شده چه چیزی را نشان می‌دهند و چه چیزی را نشان نمی‌دهند.

\textbf{روش.} یک ستون فقرات مشترک با دو سر آستانه‌ای — پنج‌درجه‌ای برای رتینوپاتی و
سه‌درجه‌ای برای ادم ماکولا — روی \textbf{\num{۲۲۶۰}} تصویر از دو پیکرهٔ عمومی آموزش دید، با
اعتبارسنجی متقاطع گروهی پنج‌تایی روی تقسیمی که پیش از نخستین اجرا منجمد و بر پایهٔ خوشه‌بندی
ادراکی ساخته شد تا نشت میان‌پیکره‌ای بسته شود. \textbf{\num{۳۶۶۲}} تصویر \lr{APTOS} تا انجماد مدل
نهایی کنار گذاشته شد. قاعدهٔ حاکم بر همهٔ مقایسه‌ها این است که \textbf{هیچ اختلافی به بازنمایی
نسبت داده نمی‌شود مگر آن‌که هر دو مدل زیر نقاط برش برازش‌متقاطع سنجیده شوند.}

\textbf{نتایج.} مدل انتخاب‌شده روی استخر توسعه به \lr{QWK} برابر \textbf{\num{۰٫۸۹۵۴}} و روی
مجموعهٔ بیرونی به \textbf{\num{۰٫۹۰۲۶}} با بازهٔ [\num{۰٫۸۹۴۱}، \num{۰٫۹۱۰۸}] می‌رسد، و خط لولهٔ کامل
\textbf{\num{۰٫۲۳۴۲}} واحد بر ویژگی‌های منجمد و یک مدل خطی می‌افزاید. از \textbf{نه} مداخلهٔ
پیش‌ثبت‌شده، \textbf{هشت} مداخله ابطال شدند. سر ادم ماکولا به سقفی می‌خورد که از داده می‌آید،
نه از معماری: شش مداخلهٔ مستقل — از جمله برشی که ناحیهٔ تعریف‌کنندهٔ درجهٔ بالینی را جدا
می‌کند — بازهٔ اطمینان را از صفر دور نکردند، و تنها \textbf{\num{۵۱}} تصویر در کل پژوهش درجهٔ
میانی دارند. یک نتیجهٔ منفی جداگانه نشان می‌دهد \textbf{درجهٔ ادم ماکولا از روی
حاشیه‌نویسی‌های منتشرشده بازسازی نمی‌شود}، پس مسیر ساختن دادهٔ سه‌کلاسهٔ بیشتر از روی
ماسک‌های قطعه‌بندی بسته است.

\textbf{سهم.} سهم اصلی این کار روش‌شناختی است. قاعدهٔ واسنجی همتاشده تا امروز چهار ادعا را در
\emph{دو جهت} برگردانده و دو بار علامت اختلاف را وارونه کرده است؛ مهم‌تر آن‌که همین قاعده کوری
مدل نسبت به درجهٔ خفیف را از یک محدودیت ظرفیت به یک مسئلهٔ واسنجی تبدیل کرد که بدون آموزش
دوباره حل می‌شود. در کنار آن، مجموعه‌ای از سازوکارهای تبار داده ارائه می‌شود که با پنج نمونه
انگیزه‌مند شده‌اند و \textbf{دو نمونه از آن‌ها شکست خودِ این پایان‌نامه است.}

\textbf{محدودیت.} تقسیم اختصاصی و منجمد این پژوهش، هر عدد این متن را با هر عدد منتشرشده‌ای
ناهم‌سنجه می‌کند. این بده‌بستان آگاهانه است: مقایسه‌پذیری فدای درستی شده.

\vspace{1em}
\noindent\textbf{کلیدواژه‌ها:} رتینوپاتی دیابتی، ادم ماکولای دیابتی، درجه‌بندی ترتیبی، واسنجی،
نقطهٔ برش، اعتبارسنجی بیرونی، تبار داده، مدل بنیادین.

\cleardoublepage

\chapter*{پیش‌گفتار}
\addcontentsline{toc}{chapter}{پیش‌گفتار}
\markboth{پیش‌گفتار}{پیش‌گفتار}

این پژوهش با هدفی آغاز شد که در میانهٔ کار عوض شد، و بهتر است خواننده این را از ابتدا بداند.

هدف نخست، ساختن مدلی با بالاترین دقت ممکن برای درجه‌بندی هم‌زمان رتینوپاتی دیابتی و ادم ماکولا
بود. آن مدل ساخته شد و نتایجش در فصل \ref{ch:results} می‌آید. ولی در مسیر ساختنش، چیزی روشن شد
که از خود مدل مهم‌تر بود: \textbf{بخش بزرگی از آنچه در این حوزه به معماری نسبت داده می‌شود،
اثر انتخاب نقاط برش است} — انتخابی که معمولاً انجام نمی‌شود بلکه از یک پیش‌فرض به ارث می‌رسد.

از آن نقطه به بعد، پرسش پژوهش عوض شد. کار به جای «کدام معماری بهتر است؟» به این پرسش پرداخت
که «چه چیزی باید ثابت بماند تا مقایسهٔ دو مدل معنا داشته باشد؟» نتیجه این بود که از نه مداخلهٔ
آزموده‌شده، هشت مداخله ابطال شدند. \textbf{این متن آن ابطال‌ها را در پیوست پنهان نمی‌کند و در
فصل \ref{ch:results} کنار هم می‌آورد}، چون فهرست چیزهایی که کار نکردند، در این پژوهش، خودِ
نتیجه است.

نکتهٔ دوم که باید پیشاپیش گفت این است که این پایان‌نامه دربارهٔ تبار داده حرف می‌زند و
\textbf{دو نمونه از نمونه‌هایش شکست خودِ ماست}. نسخه‌های نخست همین متن ادعا می‌کردند که همهٔ
عددهایشان تولیدشده است، در حالی که با دست تایپ شده بودند و سه عدد غلط بود — یکی از آن سه، شاهد
اصلی یک استدلال بود و هیچ اسکریپتی پشتش نبود. و نمونهٔ دوم هنگام نوشتن همان بخش ساخته شد: یک
تشخیص که از روی مشخصات یک استاندارد گرفته شد و نه از روی صفحهٔ چاپ‌شده، و چیزی را «تصحیح» کرد
که غلط نبود. هر دو در بخش \ref{sec:our-own-failure} گزارش شده‌اند. \textbf{آوردنشان انتخاب
است، نه اجبار}، و دلیلش این است که فصلی که تبار دادهٔ دیگران را حسابرسی کند و شکست هم‌جنس خود
را نگوید، ادعایی می‌کند که خودش را زیر آن نمی‌برد.

سرانجام، دربارهٔ آنچه این متن ادعا نمی‌کند: هیچ عددی در این پایان‌نامه با کار منتشرشده قابل
مقایسه نیست، و هیچ ادعایی دربارهٔ آمادگی بالینی در آن نیست. این یک مطالعهٔ روش‌شناختی روی
داده‌های عمومی است.

\cleardoublepage

\chapter*{فهرست نشانه‌ها و کوته‌نوشت‌ها}
\addcontentsline{toc}{chapter}{فهرست نشانه‌ها و کوته‌نوشت‌ها}
\markboth{فهرست نشانه‌ها}{فهرست نشانه‌ها}

\section*{کوته‌نوشت‌ها}

اصطلاح‌های زیر در سراسر متن به شکل لاتین می‌آیند و ترجمه نمی‌شوند، چون معادل فارسی
جاافتاده‌ای ندارند و ترجمه‌شان ارجاع به ادبیات را دشوار می‌کند.

\begin{latinlist}
\item[\lr{DR}] رتینوپاتی دیابتی \lr{(diabetic retinopathy)}
\item[\lr{DME}] ادم ماکولای دیابتی \lr{(diabetic macular edema)}
\item[\lr{QWK}] کاپای وزن‌دار درجه‌دو \lr{(quadratic weighted kappa)} — معیار اصلی این پژوهش
\item[\lr{AUROC}] سطح زیر منحنی مشخصهٔ عملکرد \lr{(area under the ROC curve)}
\item[\lr{AUPR}] سطح زیر منحنی دقت-بازخوانی \lr{(area under the precision--recall curve)}
\item[\lr{ECE}] خطای واسنجی مورد انتظار \lr{(expected calibration error)}
\item[\lr{F1}] میانگین همساز دقت و بازخوانی
\item[\lr{CFP}] عکس رنگی ته‌چشم \lr{(color fundus photograph)}
\item[\lr{OCT}] مقطع‌نگاری همدوسی نوری \lr{(optical coherence tomography)}
\item[\lr{NPDR}] رتینوپاتی دیابتی غیرتکثیری \lr{(non-proliferative DR)}
\item[\lr{PDR}] رتینوپاتی دیابتی تکثیری \lr{(proliferative DR)}
\item[\lr{TTA}] افزایش داده در زمان آزمون \lr{(test-time augmentation)}
\item[\lr{LP-FT}] کاوش خطی سپس ریزتنظیم \lr{(linear probe then fine-tune)}
\item[\lr{OOF}] برون‌تایی \lr{(out-of-fold)}
\item[\lr{DD}] قطر دیسک نوری \lr{(disc diameter)} — واحد فاصله در تعریف بالینی درجهٔ ادم ماکولا
\item[\lr{GPU}] پردازندهٔ گرافیکی
\end{latinlist}

\section*{نشانه‌های ریاضی}

\begin{latinlist}
\item[$K$] شمار درجه‌ها در یک سر؛ $K=5$ برای رتینوپاتی و $K=3$ برای ادم ماکولا
\item[$y$] درجهٔ مرجع یک تصویر، $y \in \{0,\dots,K-1\}$
\item[$P(y>k)$] احتمالی که سر آستانه‌ای برای آستانهٔ $k$ تخمین می‌زند
\item[$s$] نمرهٔ پیوسته، $s=\sum_{k=0}^{K-2} P(y>k)$، یعنی درجهٔ مورد انتظار
\item[$c_k$] نقطهٔ برش $k$-ام؛ رمزگشایی، $s$ را با این نقاط به یک درجه تبدیل می‌کند
\item[$\kappa$] کاپای وزن‌دار درجه‌دو
\item[$\alpha,\beta$] وزن دو سر در تابع زیان مشترک
\end{latinlist}

\section*{قراردادهای نگارشی عددها}

عددهای متن با \textbf{ارقام فارسی} و جداکنندهٔ اعشار «٫» می‌آیند؛ عددهای درون فرمول با ارقام
لاتین. \textbf{هر عدد فارسی این متن از پروندهٔ تولیدشدهٔ
\lr{docs/generated/thesis\_numbers.json} می‌آید} و یک آزمون خودکار هر رقمی را که در آن پرونده
نباشد رد می‌کند. دلیل وجود این سازوکار در بخش \ref{sec:our-own-failure} آمده است.

\cleardoublepage

% English abstract — set left-to-right inside the Persian document.

\chapter*{Abstract}
\addcontentsline{toc}{chapter}{Abstract (English)}
\markboth{Abstract}{Abstract}
\begin{latin}
\noindent
Diabetic retinopathy is the leading cause of preventable blindness among working-age adults, and
diabetic macular edema is the most common mechanism of central vision loss. Automated fundus
grading reports high headline figures, yet as of the end of April 2024 only three ophthalmic AI
devices had United States Food and Drug Administration clearance. This thesis begins from that
gap and asks what the reported numbers do and do not establish.

\vspace{0.6em}
\noindent\textbf{Methods.} A shared backbone with two threshold heads --- five ordinal grades for
retinopathy and three for macular edema --- was trained on 2{,}260 images from two public corpora
under five-fold grouped cross-validation. The split was frozen before the first real run and
built from perceptual clustering so that cross-corpus leakage is closed; 3{,}662 APTOS images were
held out until the final model was frozen. One rule governs every comparison in this work:
\textbf{no difference is attributed to representation unless both models are evaluated under
cross-fitted cut-points}. An untuned default is a hyper-parameter that has been chosen, not one
that has been avoided.

\vspace{0.6em}
\noindent\textbf{Results.} The selected model reaches a quadratic weighted kappa of 0.8954 on the
development pool and 0.9026 [0.8941, 0.9108] on the held-out external set, and the full pipeline
is worth +0.2342 kappa over frozen ImageNet features with a linear ordinal head on identical
images, folds, decoding and calibration. Of nine pre-registered interventions, eight were
falsified. The macular-edema head meets a ceiling that comes from data rather than architecture:
six independent interventions --- including a crop isolating the region the clinical grade is
defined on --- failed to move the confidence interval away from zero, and only 51 images in the
entire study carry the middle grade. A separate negative result establishes that the macular-edema
grade \textbf{cannot be recovered from published annotations}, closing the route to more
three-class data via segmentation masks.

\vspace{0.6em}
\noindent\textbf{Contribution.} The contribution is methodological rather than a performance
figure. The matched-calibration rule has reversed four prior attributions \emph{in both
directions}, with two sign reversals; more importantly, the same rule converted the model's
blindness to the mild grade from a capacity limit into a calibration problem solvable without
retraining. Alongside it, a set of data-provenance mechanisms is presented, motivated by five
worked examples --- \textbf{two of which are this thesis's own failures}: a chapter that
asserted its numbers were generated while they were typed by hand, three of them wrong and one
supporting a central argument with no artefact behind it at all; and, committed while writing
that up, a diagnosis taken from a specification rather than from the rendered page, which
``corrected'' something that was not broken.

\vspace{0.6em}
\noindent\textbf{Limitation.} The frozen, project-specific split makes every number here
incommensurable with any published figure. That trade is deliberate: comparability was given up
for correctness, and the reader should know which of the two they hold.

\vspace{1em}
\noindent\textbf{Keywords:} diabetic retinopathy, diabetic macular edema, ordinal grading,
calibration, decision threshold, external validation, data provenance, foundation models.
\end{latin}
\cleardoublepage


\tableofcontents
\cleardoublepage
\listoftables
\cleardoublepage

\cleardoublepage
\pagenumbering{arabic}

% ------------------------------------------------------------------ فصل‌ها
% فصل ۱ — مقدمه
% همهٔ عددها از docs/generated/thesis_numbers.json می‌آیند، که src/thesis_numbers.py آن را
% از آرشیو runs/ و docs/generated/ می‌سازد. صحت با «python3 tools/farsi_lint.py --numbers» بررسی می‌شود.

\chapter{مقدمه}
\label{ch:intro}

\section{مسئله}

رتینوپاتی دیابتی بزرگ‌ترین علت نابینایی پیشگیری‌پذیر در جمعیت شاغل جهان است، و ادم ماکولای
دیابتی شایع‌ترین سازوکاری است که در آن بینایی مرکزی از دست می‌رود. هر دو با غربالگری منظم و
درمان به‌موقع قابل کنترل‌اند، و هر دو تا مرحله‌ای پیش می‌روند که بیمار علامتی حس نکند. پس
مسئله، پیش از آن‌که مسئلهٔ درمان باشد، مسئلهٔ \emph{ظرفیت غربالگری} است: شمار افراد دیابتی از
شمار متخصصانی که می‌توانند ته‌چشمشان را درجه‌بندی کنند بسیار بیشتر است.

همین شکاف، دلیل سه دههٔ کار روی درجه‌بندی خودکار تصویر ته‌چشم است. مدل‌های یادگیری عمیق روی
این وظیفه اعدادی گزارش می‌کنند که در نگاه نخست مسئله را حل‌شده نشان می‌دهد. این پایان‌نامه از
همان‌جا آغاز می‌کند که آن اعداد چه چیزی را نشان می‌دهند و چه چیزی را نشان نمی‌دهند.

\section{شکافی که بسته نمی‌شود}

اگر مسئله حل شده بود، سامانه‌های خودکار در بالین می‌بودند. نیستند. تا پایان آوریل \num{۲۰۲۴} تنها
\textbf{\num{۳}} دستگاه هوش مصنوعی چشم‌پزشکی از سازمان غذا و داروی ایالات متحده مجوز گرفته
بودند، و هر سه تنها برای غربالگری رتینوپاتی دیابتی و تنها با دوربین‌های مشخص
\cite{ruamviboonsuk2024review}. مرور یادشده نتیجه را صریح می‌نویسد: شمار پژوهش‌های روش‌های
تازه نمایی بالا می‌رود، در حالی که پذیرش در مراقبت چشمی با شیبی بسیار کمتر رشد می‌کند، پس
\textbf{شکاف میان توسعه و استقرار در حال گشاد شدن است}.

همان مرور چهار مانع را نام می‌برد، و مانع نخست انگیزهٔ این پایان‌نامه است: \textbf{کمبود
مقایسه‌های سربه‌سر میان مدل‌های موجود}. سه مانع دیگر — نبود شواهد روشن هزینه‌اثربخشی، مسائل
عدالت و سوگیری، و مسئولیت حقوقی-پزشکی — بیرون از دامنهٔ این کارند و در فصل \ref{ch:discussion}
به‌عنوان محدودیت نام برده می‌شوند.

مرور نظام‌مند دیگری، روی \textbf{\num{۳۸}} پژوهش، نشان می‌دهد این کمبود از کجا می‌آید: تنها
\textbf{\num{۱۲}} پژوهش، یعنی \textbf{\num{۳۲}} درصد، اعتبارسنجی بیرونی انجام داده‌اند، و تنها
\textbf{\num{۴}} پژوهش، یعنی \textbf{\num{۱۱}} درصد، شیوهٔ برخورد با داده‌های ازدست‌رفته را گزارش
کرده‌اند \cite{yang2025review}. همان مرور، در بخش بحث، نتیجه‌ای می‌گیرد که در سراسر این
پایان‌نامه تکرار خواهد شد: مدل‌هایی که تنها روی دادهٔ درونی آموزش دیده‌اند کارایی بالاتری
گزارش می‌کنند، که نشانهٔ خطر بیش‌برازش و سوگیری آموزش است.

\section{آنچه یک عدد سرتیتر پنهان می‌کند}

بین «مدل خوب کار می‌کند» و «مدل در بالین کار می‌کند» یک لایهٔ نادیده وجود دارد که این
پایان‌نامه دربارهٔ آن است: \textbf{قاعدهٔ تصمیم}.

یک مدل درجه‌بندی، نمره‌ای پیوسته تولید می‌کند. آنچه بالین می‌بیند یک درجه است، و میان آن دو،
مجموعه‌ای از نقاط برش می‌نشیند. این نقاط تقریباً همیشه انتخاب نمی‌شوند: آستانهٔ پیش‌فرض
\wrongnum{۵/۰} روی خروجی سیگموئید — که در این متن به‌درستی \num{۰٫۵} نوشته می‌شود — از مقیاس
درجه‌بندی نمی‌آید، از تابع فعال‌سازی می‌آید. \textbf{یک پیش‌فرض تنظیم‌نشده، فراپارامتری است
که انتخاب شده، نه فراپارامتری که از آن پرهیز شده است.}

اهمیت این نکته را ادبیات موجود، بی‌آن‌که نتیجه‌اش را بگیرد، بارها نشان داده است. در یک
اعتبارسنجی واقعی از سامانه‌های تجاری، ویژگی سه فروشنده روی \emph{همان چشم‌ها} و با همان مرجع،
از \textbf{\num{۱۴٫۲۵}} درصد تا \textbf{\num{۹۶٫۰۱}} درصد نوسان می‌کند \cite{duggal2025realworld}. در یک
فراتحلیل، حساسیت تجمیعی \textbf{\num{۰٫۹۵}} گزارش می‌شود در حالی که پژوهش‌های واردشده بازه‌ای تا
\textbf{\num{۰٫۰۶}} را در بر می‌گیرند \cite{tahir2025meta}. و در یک محک ترتیبی، همان قاعدهٔ آستانه
که کم‌برآوردی درجه‌های شدید را از \textbf{\num{۸۴٫۷}} به \textbf{\num{۴۹٫۷}} می‌رساند، شمار
پیش‌بینی‌های درجهٔ یک را به \textbf{\num{۷۱۲٫۷}} در برابر \textbf{\num{۲۷۰}} مورد واقعی می‌رساند، با
دقت \textbf{\num{۰٫۱۹۳}} \cite{sheng2026orderdr}.

\textbf{هیچ‌یک از این مقاله‌ها یک جاروب آستانه گزارش نمی‌کند.} همه اثر آن را می‌بینند و
هیچ‌کدام آن را به‌عنوان متغیر مطالعه نمی‌کنند.

\section{پرسش‌های این پژوهش}

از این‌جا سه پرسش می‌آید که ساختار پایان‌نامه را می‌سازند.

\textbf{پرسش یکم — قاعدهٔ تصمیم چقدر از آنچه به معماری نسبت می‌دهیم را توضیح می‌دهد؟} اگر
جابه‌جایی نقاط برش بتواند اختلاف میان دو معماری را از میان ببرد یا علامتش را عوض کند، آنگاه
هر مقایسهٔ معماری که نقاط برش را همتا نکرده باشد، دربارهٔ آستانه حرف می‌زند و فکر می‌کند
دربارهٔ بازنمایی حرف می‌زند.

\textbf{پرسش دوم — سقف کارایی سر ادم ماکولا کجاست، و از داده می‌آید یا از معماری؟} این پرسش
به این دلیل قابل پاسخ است که می‌توان مداخله‌هایی طراحی کرد که مستقیماً ناحیهٔ تعریف‌کنندهٔ
برچسب را هدف بگیرند؛ اگر آن‌ها هم کاری نکنند، محدودیت جای دیگری است.

\textbf{پرسش سوم — تبار داده در این حوزه چقدر قابل اتکاست، و برای اتکاپذیر کردنش چه چیزی
لازم است؟} این پرسش در آغاز کار مطرح نبود و از دل کار بیرون آمد.

\section{سهم این پایان‌نامه}

سهم اصلی این کار \textbf{روش‌شناختی} است، نه کارایی. این را باید در آغاز نوشت، چون خواننده‌ای
که به دنبال عددی بهتر آمده است باید بداند کجا را نگاه کند و کجا را نه.

\textbf{یکم، قاعدهٔ واسنجی همتاشده و سابقهٔ عمل آن.} هیچ اختلافی به بازنمایی نسبت داده نمی‌شود
مگر آن‌که هر دو مدل زیر نقاط برش برازش‌متقاطع سنجیده شوند. این قاعده تا امروز چهار ادعا را در
\emph{دو جهت} برگردانده و دو بار علامت اختلاف را وارونه کرده است، و سابقه‌اش در فصل
\ref{ch:results} می‌آید. مهم‌تر از هر تک‌نتیجه، این است که همان قاعده یک نتیجهٔ ظاهراً منفی —
کوری نسبت به درجهٔ خفیف — را به مسئله‌ای بی‌هزینه و حل‌شدنی تبدیل کرد.

\textbf{دوم، یک سقف نظارت اندازه‌گیری‌شده برای ادم ماکولا.} شش مداخلهٔ مستقل، از جمله
مداخله‌ای که مستقیماً ناحیهٔ تعریف‌کنندهٔ درجهٔ بالینی را جدا می‌کند، بازهٔ اطمینان را از صفر
دور نکردند. این نتیجه با شاهدی بیرونی تقویت می‌شود که همهٔ توضیح‌های جایگزین را می‌بندد: یک
پیکرهٔ ادم ماکولای ساخته‌شده برای همین کار، با \textbf{\num{۲۹۳۸}} تصویر، باز هم تنها
\textbf{\num{۹٫۵}} درصد نمونه در درجهٔ میانی دارد \cite{tang2026mmrdr}.

\textbf{سوم، تبار داده به‌عنوان سهم، نه به‌عنوان اعتراف.} شش اشکال تبار داده در این پژوهش هیچ
خطای زمان اجرا تولید نکردند و در سکوت عدد غلط ساختند. سازوکارهایی که پاسخ آن‌هاست در فصل
\ref{ch:method} می‌آید، و \textbf{یکی از پنج نمونهٔ آن فصل، شکست خودِ این پایان‌نامه است}.

\textbf{چهارم، یک نتیجهٔ منفی که راهی را می‌بندد.} درجهٔ ادم ماکولا با تعریف هندسی منتشرشده‌اش
از روی حاشیه‌نویسی‌های منتشرشده بازسازی نمی‌شود. این را فصل \ref{sec:parta} نشان می‌دهد و
پیامدش این است که مسیر ساختن دادهٔ سه‌کلاسهٔ بیشتر از روی ماسک‌های قطعه‌بندی بسته است.

\section{آنچه این پایان‌نامه ادعا نمی‌کند}

\textbf{هیچ عددی در این متن با کار منتشرشده قابل مقایسه نیست.} تقسیم این پژوهش اختصاصی و
منجمد است و بر پایهٔ خوشه‌بندی ادراکی ساخته شده تا نشت میان‌پیکره‌ای را ببندد؛ همین کار آن را
با هر تقسیم منتشرشده‌ای ناهم‌سنجه می‌کند. این محدودیت در فصل \ref{ch:results} و فصل
\ref{ch:discussion} تکرار می‌شود، چون تنها جایی که می‌تواند آسیب بزند همان جایی است که
خواننده عادت دارد عددها را کنار هم بگذارد.

\textbf{و ادعای برتری معماری هم در کار نیست.} از نه مداخلهٔ آزموده‌شده، هشت مداخله ابطال شدند.
آنچه می‌ماند رویه‌ای است برای رسیدن به عددی که بتوان از آن دفاع کرد.

\section{ساختار پایان‌نامه}

فصل \ref{ch:related} ادبیات را نه بر پایهٔ موضوع، که بر پایهٔ \emph{تناقض‌ها} سازمان می‌دهد:
سه جایی که نتیجه‌های منتشرشده با هم نمی‌خوانند، و این‌که هر تناقض در کدام خانهٔ آزموده‌نشده
حل می‌شود. فصل \ref{ch:method} پروتکل ارزیابی، داده، معماری و سازوکارهای تبار داده را می‌آورد؛
معماری آخرین چیزی است که در آن فصل اهمیت پیدا می‌کند. فصل \ref{ch:results} نتایج را می‌آورد،
با ترتیبی که از «خط لوله در برابر چه چیزی سنجیده می‌شود» آغاز می‌کند. فصل \ref{sec:parta}
گزارش مستقل نتیجهٔ منفی بازسازی درجه است. فصل \ref{ch:discussion} بحث، محدودیت‌ها و کارهای
آینده را جمع می‌بندد.

% فصل ۲ — مرور ادبیات
% همهٔ عددها از docs/generated/thesis_numbers.json می‌آیند و در results/VERIFICATION-*.md
% در برابر منبع اصلی راستی‌آزمایی شده‌اند. ادعای راستی‌آزمایی‌نشده در این فصل نقل نمی‌شود.

\chapter{مرور ادبیات}
\label{ch:related}

این فصل ادبیات را بر پایهٔ موضوع سازمان نمی‌دهد، بلکه بر پایهٔ \textbf{تناقض} سازمان می‌دهد.
دلیلش این است که فهرستی از «چه کسی چه کرد» به پرسش‌های فصل \ref{ch:intro} پاسخ نمی‌دهد، در
حالی که جاهایی که نتیجه‌های منتشرشده با هم نمی‌خوانند دقیقاً همان‌جاهایی‌اند که یک آزمایش تازه
چیزی به دست می‌آورد.

سه تناقض بررسی می‌شود. در هر سه، الگو یکسان است: \textbf{نتیجه‌های ناسازگار در خانه‌هایی
نشسته‌اند که مقالهٔ مبدأ آن‌ها را نیازموده است.} پس تناقض‌ها با انتخاب یک طرف حل نمی‌شوند، با
نام بردن خانهٔ آزموده‌نشده حل می‌شوند.

\paragraph{یادداشتی دربارهٔ اتکاپذیری این فصل.} هر ادعای نقل‌شده در این فصل در برابر متن اصلی
مقاله بررسی شده است و پرونده‌های راستی‌آزمایی در مخزن پژوهش نگهداری می‌شوند. جایی که عددی در
منبع پیدا نشد یا با آنچه انتظار می‌رفت نخواند، \textbf{آن عدد در این متن نقل نمی‌شود}. سه مورد
از این دست وجود دارد و در بخش \ref{sec:lit-corrections} نام برده می‌شوند، چون خودِ آن‌ها
یافته‌اند.

\section{مدلی که زیر حسابرسی است}
\label{sec:retfound-audit}

نقطهٔ مرجع این ادبیات یک مدل بنیادین شبکیه است که با یادگیری خودنظارتی روی \textbf{\num{۱٫۶}}
میلیون تصویر ته‌چشم و مقطع‌نگاری همدوسی نوری پیش‌آموزش دیده و برای وظایف پایین‌دستی ریزتنظیم
می‌شود \cite{zhou2023retfound}. مروری داوری‌شده آن را چنین خلاصه می‌کند: این مدل خودنظارتی
مدل‌های نظارت‌شده و خودنظارتیِ پیش‌آموزش‌دیده روی \lr{ImageNet} را در همان وظایف، \emph{هم در
اعتبارسنجی درونی و هم بیرونی}، پشت سر می‌گذارد \cite{ruamviboonsuk2024review}.

\textbf{این ادعا زیر حسابرسی است، نه زیر سؤال.} تفاوت مهم است. مقالهٔ مبدأ آنچه را آزموده
درست گزارش می‌کند؛ آنچه این فصل نشان می‌دهد این است که \emph{بخش بزرگی از فضای مسئله را
نیازموده}، و هر نتیجهٔ ناسازگارِ بعدی دقیقاً در همان بخش نشسته است. مشخصاً: در آن مقاله هیچ
پایهٔ شبکهٔ همگشتی وجود ندارد، هیچ رمزگذار منجمد یا کاوشگر خطی ارزیابی نمی‌شود، هیچ معیار
به‌تفکیک‌کلاس برای رتینوپاتی گزارش نمی‌شود، و واسنجی زیر جابه‌جایی توزیع سنجیده نمی‌شود.

\section{تناقض یکم: مدل بنیادین منجمد در برابر ریزتنظیم‌شده}
\label{sec:contradiction1}

\subsection{شواهدی که با هم نمی‌خوانند}

سه گروه مستقل، سه پاسخ متفاوت به «آیا مدل بنیادین شبکیه بهتر است؟» می‌دهند.

\textbf{پاسخ یکم: بله، و تفاوت ناچیز است.} روی یک پیکرهٔ چندپیمانه‌ای، دقت درجه‌بندی
رتینوپاتی روی تصویر ته‌چشم برای مدل بنیادین \textbf{\num{۰٫۸۲۲}} و برای یک شبکهٔ همگشتی
پیش‌آموزش‌دیده روی \lr{ImageNet} برابر \textbf{\num{۰٫۸۲۳}} است — عملاً مساوی
\cite{tang2026mmrdr}.

\textbf{پاسخ دوم: بله، ولی فقط درون توزیع.} یک سر لجستیک با \textbf{\num{۵۱۲۵}} پارامتر روی
ویژگی‌های منجمد همان مدل بنیادین، با ریزتنظیم کامل همان مدل که \textbf{\num{۳۰۳٫۳}} میلیون
پارامتر دارد تفکیک‌ناپذیر است \cite{nielsen2025he}. اعداد این‌ها هستند: روی
\lr{APTOS} سر منجمد \textbf{\num{۰٫۹۳}} در برابر ریزتنظیم \textbf{\num{۰٫۹۴}}، و روی \lr{ODIR-5K}
\textbf{\num{۰٫۷۸}} در برابر \textbf{\num{۰٫۸۰}}. \textbf{در هر دو مورد سر منجمد اسماً عقب‌تر است و
اختلاف معنادار نیست.} نویسندگان خود می‌نویسند که ارزیابی بیرون از توزیع را انجام نداده‌اند و
انجامش بینش ارزشمندی می‌داد.

\textbf{پاسخ سوم: نه، وقتی پیمانه عوض شود.} روی همان پیکرهٔ چندپیمانه‌ای، هنگام انتقال به
تصویربرداری فوق‌عریض، شبکهٔ همگشتی پیش‌آموزش‌دیده روی \lr{ImageNet} \emph{بهتر} از مدل‌های
بنیادین چشم‌پزشکی تعمیم می‌دهد \cite{tang2026mmrdr}.

\subsection{چرا این‌ها در واقع تناقض نیستند}

هر سه نتیجه در خانه‌ای نشسته‌اند که مقالهٔ مبدأ نیازموده است. مقالهٔ مبدأ رمزگذار منجمد را
ارزیابی نمی‌کند، پس با پاسخ دوم برخورد نمی‌کند. پایهٔ همگشتی ندارد، پس با پاسخ‌های یکم و سوم
برخورد نمی‌کند. \textbf{تناقض ظاهری، محصول تعمیم دادن نتیجه‌ای است به شرایطی که در آن آزموده
نشده.}

نتیجهٔ خود این پژوهش در فصل \ref{ch:results} در همین شکاف می‌نشیند و آن را تیزتر می‌کند: با
رمزگذار منجمد، بازنمایی مدل بنیادین \emph{بهتر} است؛ با ریزتنظیم کامل، همان مدل به خط لولهٔ
ساده‌تر ما \emph{می‌بازد}. دو سازوکار نامزد وجود دارد و شواهد ما آن‌ها را از هم جدا نمی‌کند؛
این در فصل \ref{ch:discussion} صریح گفته می‌شود.

\subsection{یک اشکال روش‌شناختی که در همین‌جا تکرار می‌شود}

پیکرهٔ چندپیمانه‌ای یادشده مدل‌ها را کنار هم جدول می‌کند در حالی که \emph{شیوهٔ تطبیق آن‌ها
یکسان نیست}: برخی به‌طور کامل ریزتنظیم می‌شوند و برخی با «تنظیم حداقلی» ارزیابی می‌شوند. این
دقیقاً همان مخدوش‌کننده‌ای است که فصل \ref{ch:method} در بخش \ref{sec:matched} می‌بندد.

\textbf{و نمونهٔ خلافش در ادبیات وجود دارد.} مقالهٔ مبدأ، هنگام مقایسهٔ راهبردهای پیش‌آموزش،
معماری و فرایند تطبیق را ثابت نگه می‌دارد و تنها چیز تحت آزمون را تغییر می‌دهد، و خود هشدار
می‌دهد که ادعای برتری یک روش با وجود چند متغیر هم‌زمان نیازمند احتیاط است
\cite{zhou2023retfound}. \textbf{این طراحی درست است و باید معیار سنجش بقیه باشد.}

\section{تناقض دوم: تشخیص‌پذیری به‌جای واسنجی}
\label{sec:contradiction2}

\subsection{خطا در طبیعت}

یک مقالهٔ داوری‌شدهٔ \num{۲۰۲۶}، پس از گزارش مقادیر \lr{AUC}، می‌نویسد که این مقادیر بالا «تأیید
می‌کنند» مدل نه‌تنها دقیق است بلکه احتمال‌های خوش‌واسنجی نسبت می‌دهد و برای پشتیبانی تصمیم
بالینی قابل اتکاست \cite{yu2026dualswinord}.

\textbf{این گزاره نادرست است، و نادرستی‌اش ساختاری است نه کمّی.} \lr{AUC} تنها به
\emph{ترتیب} نمره‌ها حساس است و نسبت به هر تبدیل یکنوا ناوردا است. مدلی که همهٔ احتمال‌هایش را
در عدد ثابتی ضرب کند، \lr{AUC} یکسانی دارد و واسنجی کاملاً متفاوتی. آن مقاله \emph{هیچ} معیار
واسنجی گزارش نمی‌کند — نه \lr{ECE}، نه \lr{Brier}، نه منحنی اطمینان — و هیچ اعتبارسنجی بیرونی
هم ندارد: هر دو محکش، پیکره‌هایی‌اند که روی آن‌ها آموزش دیده است.

\subsection{شاهد نقض، از دل همان ادبیات}

مسئله را نمی‌توان با شهود رد کرد؛ با اندازه‌گیری رد می‌شود، و اندازه‌گیری‌اش منتشر شده است. یک
محک واسنجی‌محور نشان می‌دهد که همان مدل بنیادین، \emph{کمترین} خطای واسنجی بیرونی را در کنار
\emph{بدترین} تشخیص‌پذیری بیرونی به دست می‌آورد \cite{poyrazer2026calibration}. دو معیار در
جهت مخالف هم حرکت می‌کنند، پس هیچ‌کدام جانشین دیگری نیست.

\textbf{و واسنجی زیر جابه‌جایی توزیع فرو می‌ریزد.} در یک محک ترتیبی، خطای واسنجی از
\textbf{\num{۰٫۰۴۹}} روی مجموعهٔ درونی به \textbf{\num{۰٫۱۶۰}} روی مجموعهٔ بیرونی می‌رسد
\cite{sheng2026orderdr}. سه برابر شدن خطای واسنجی، بدون آن‌که مدل عوض شود.

\section{تناقض سوم: ادعاهای معماری که برمی‌گردند}
\label{sec:contradiction3}

یک مرور داوری‌شده می‌نویسد پژوهش‌های بسیاری برتری ترانسفورمر بینایی را بر شبکهٔ همگشتی برای
وظایف رایج، از جمله غربالگری رتینوپاتی، یافته‌اند. \emph{در همان مرور} و بدون آشتی دادن این
دو، مقایسه‌ای سربه‌سر میان هشت مدل همگشتی و نه مدل ترانسفورمر گزارش می‌شود که در آن
\textbf{همهٔ} مدل‌های همگشتی از \textbf{همهٔ} مدل‌های ترانسفورمر بهترند: حساسیت و ویژگی
مدل‌های همگشتی به \num{۹۰} درصد یا بیشتر می‌رسد، در حالی که ترانسفورمرها به حساسیت \num{۶۳} تا \num{۹۴} درصد و
ویژگی \num{۲۴} تا \num{۴۸} درصد می‌رسند \cite{ruamviboonsuk2024review}.

\textbf{این تناقض در همان صفحه است و آشتی داده نمی‌شود.} خواندن درست آن این نیست که یکی از دو
گزاره غلط است؛ این است که \textbf{ادعای معماری، وقتی شرایط ارزیابی عوض شود، برمی‌گردد}. نتیجهٔ
خود این پژوهش در فصل \ref{ch:results} همین را به شکل کمّی می‌گوید: از نه مداخله، هشت مداخله
ابطال شدند، و در یک مورد \emph{علامت} اختلاف با جابه‌جایی نقاط برش وارونه شد.

\section{خوشهٔ دوم: شواهد استقرار}
\label{sec:deployment}

خوشهٔ دوم ادبیات، مقالهٔ توسعهٔ مدل نیست؛ شواهد میدانی است. اهمیتش این است که استدلال قاعدهٔ
تصمیم را از یک نکتهٔ روش‌شناختی به یک مسئلهٔ بالینی تبدیل می‌کند.

\textbf{نقطهٔ عمل، نه مدل، خروجی را تعیین می‌کند.} در اعتبارسنجی واقعی سه سامانهٔ تجاری روی
همان چشم‌ها و با همان مرجع، ویژگی از \textbf{\num{۱۴٫۲۵}} تا \textbf{\num{۹۶٫۰۱}} درصد نوسان می‌کند
\cite{duggal2025realworld}. در همان کار، شاخص توافق از \textbf{\num{۰٫۶۵}} به \textbf{\num{۰٫۷۲}}
می‌رود در حالی که حساسیت و ویژگی هر کدام حدود سی‌ویک واحد جابه‌جا می‌شوند.

\paragraph{و این جابه‌جایی را نباید جاروب آستانه خواند.} دو مقدار توافق از دو فاز متفاوت، دو
گروه بیمار متفاوت و پس از بازآموزی الگوریتم توسط فروشنده می‌آیند — خودِ مقاله این را می‌گوید.
\textbf{پس آن را به‌عنوان شاهدِ بی‌حسی آماره‌های توافق تجمیعی نسبت به تغییر نقطهٔ عمل نقل
می‌کنیم، نه به‌عنوان جاروب آستانه روی دادهٔ ثابت.} این تمایز را در جای دیگری از این متن نباید
از دست داد.

\textbf{مرز خطا جای مشخصی دارد.} در پژوهشی روی \textbf{\num{۳۲۰}} چشم از \textbf{\num{۱۶۰}} بیمار، سه
روش غربالگری در تشخیص درجه‌های متوسط و شدید و تکثیری هیچ تفاوتی ندارند، و تمام اختلافشان در
چشم‌های بدون رتینوپاتی و چشم‌های با نشانه‌های خفیف است \cite{tomic2024handheld}. توافق
\textbf{\num{۰٫۸۶}} در برابر معاینهٔ بالینی و \textbf{\num{۰٫۷۸}} در برابر دوربین استاندارد است، و
اختلاف‌ها \textbf{\num{۲۲}} و \textbf{\num{۳۴}} چشم‌اند.

\textbf{و سر ادم ماکولا در بالین همان‌جایی می‌شکند که در آزمایشگاه.} یک سامانهٔ تجاری مستقر،
در یک اجرا و روی همان تصویرها، برای رتینوپاتی نیازمند ارجاع به توافق \textbf{\num{۰٫۸۱}} می‌رسد و
برای ادم ماکولا روی تصویر ته‌چشم به حساسیت \textbf{\num{۲۶٫۵}} درصد و توافق \textbf{\num{۰٫۳۸}}
\cite{duggal2025realworld}. \textbf{این تک‌شاهد، همهٔ توضیح‌های جایگزین برای سقف ادم ماکولا را
می‌بندد}: نه معماری ما، نه اندازهٔ دادهٔ ما، نه بودجهٔ آموزش ما، نه انتخاب پیکره.

\section{شکاف تبار داده}
\label{sec:lit-provenance}

آخرین موضوع این فصل، موضوعی است که ادبیات به‌ندرت دربارهٔ خودش گزارش می‌کند.

یک فراتحلیل، حساسیت تجمیعی \textbf{\num{۰٫۹۵}} گزارش می‌کند در حالی که پژوهش‌های واردشده بازه‌ای
تا \textbf{\num{۰٫۰۶}} را در بر می‌گیرند \cite{tahir2025meta}؛ و بررسی جدول‌های همان کار نشان
می‌دهد ستون منفی‌های کاذب رونویسی ستون مثبت‌های کاذب است، و همان خرابی به حساسیت‌های
گزارش‌شده نشت کرده است. یک توصیف‌نامهٔ داده در خانوادهٔ \lr{Nature} پیکرهٔ خود را از یک
مجموعهٔ عمومی می‌سازد که در پیش‌آموزش یکی از مدل‌های محک‌شده به کار رفته، بی‌آن‌که این
هم‌پوشانی را جایی بیاورد \cite{tang2026mmrdr}؛ و زیرمجموعه‌های خود را در دو پیمانه در سطح
بیمار تقسیم می‌کند «برای جلوگیری از نشت داده» و در سومی در سطح تصویر، چون پیکرهٔ مبدأ شناسهٔ
بیمار ندارد. یک پژوهش بالینی، \textbf{\num{۳۲۰}} چشم از \textbf{\num{۱۶۰}} بیمار را چشم‌به‌چشم تحلیل
می‌کند بی‌آن‌که برای همبستگی دو چشم یک نفر تصحیحی بگذارد \cite{tomic2024handheld}.

\textbf{هیچ‌کدام از این‌ها خطای زمان اجرا تولید نمی‌کند.} همه در سکوت عدد می‌سازند. فصل
\ref{ch:method} سازوکارهایی را می‌آورد که در این پژوهش پاسخ آن‌هاست، و در بخش
\ref{sec:our-own-failure} \textbf{شکست هم‌جنس خودِ این پایان‌نامه} را در کنار همین چهار نمونه
می‌گذارد.

\section{سه تصحیح که خودشان یافته‌اند}
\label{sec:lit-corrections}

راستی‌آزمایی این فصل در برابر متن اصلی، سه موردی را پیدا کرد که نقلشان بدون تصحیح، خطا را
تکثیر می‌کرد. ثبتشان اینجا آمده چون \emph{نمونهٔ همان مسئله‌ای‌اند که این فصل دربارهٔ آن است}.

\textbf{یکم، معیار عوض می‌شود.} مقادیر \textbf{\num{۰٫۸۲۲}} و \textbf{\num{۰٫۸۲۳}} در بخش
\ref{sec:contradiction1} \emph{دقت}‌اند، نه \lr{AUC}. آن مقاله صریحاً می‌نویسد که \lr{AUC} را
به دلیل هزینهٔ محاسباتی کنار گذاشته است. نقل کردنشان به‌عنوان \lr{AUC} معیاری را به مقاله
نسبت می‌داد که خودش از محاسبه‌اش سر باز زده است.

\textbf{دوم، جهت عوض می‌شود.} در مقایسهٔ سر منجمد با ریزتنظیم، سر منجمد روی هر دو مجموعه
اسماً \emph{عقب‌تر} است، نه هم‌تراز یا جلوتر. نتیجهٔ «تفکیک‌ناپذیر» سرجایش می‌ماند، ولی جهتی
که خلاصه‌های ثانویه القا می‌کنند در منبع وجود ندارد.

\textbf{سوم، نرخی که بر پایهٔ سال اشتباه بنا شده.} مرور یادشده پس از گزارش درست
\textbf{\num{۳}} دستگاه دارای مجوز، نرخ «یک مدل در سال» را از این فرض می‌گیرد که نخستین مجوز در
\num{۲۰۲۱} صادر شده است؛ نخستین مجوز در \num{۲۰۱۸} صادر شد. \textbf{شمار دستگاه‌ها نقل می‌شود، نرخ نه.}

یک ادعای چهارم هم هست که در این پایان‌نامه \emph{اصلاً} نقل نمی‌شود، چون در متن اصلی یافت نشد
و متن اصلی ظاهراً خلافش را می‌گوید. تصمیم بر این شد که به‌جای بازنویسی محتاطانه، کنار گذاشته
شود: \textbf{ادعایی که راستی‌آزمایی نشده، نقل نمی‌شود.}

\section{جای این پژوهش}

سه تناقض این فصل، سه خانهٔ آزموده‌نشده را نام می‌برند: مقایسهٔ منجمد و ریزتنظیم‌شده در یک
طراحی که تنها یک چیز را عوض کند؛ واسنجی به‌عنوان معیاری مستقل از تشخیص‌پذیری و زیر جابه‌جایی
توزیع؛ و ادعای معماری زیر نقاط برش همتاشده. فصل \ref{ch:method} پروتکلی را می‌سازد که هر سه
را هم‌زمان می‌بندد، و فصل \ref{ch:results} نشان می‌دهد وقتی این کار انجام شود، بیشتر
مداخله‌ها چیزی نمی‌سازند — و همین، نتیجه است.

% فصل ۳ — روش کار
% همهٔ عددها از docs/generated/thesis_numbers.json می‌آیند، که src/thesis_numbers.py آن را
% از آرشیو runs/ و docs/generated/ می‌سازد. صحت با «python3 tools/farsi_lint.py --numbers» بررسی می‌شود.

\chapter{روش کار}
\label{ch:method}

این فصل تصمیم‌های روشی را توصیف می‌کند. سهم اصلی این پژوهش روش‌شناختی است، نه کارایی؛ پس
ترتیب فصل نیز همین را دنبال می‌کند: نخست پروتکل ارزیابی، بعد داده، بعد معماری، و در پایان
سازوکارهایی که تبار داده را تضمین می‌کنند. معماری آخرین چیزی است که اهمیت پیدا می‌کند، و
فصل \ref{ch:results} نشان می‌دهد چرا.

\section{پروتکل ارزیابی}

پروتکل پیش از نخستین اجرای واقعی نوشته و منجمد شد. تغییر آن نیازمند افزودن بخشی تاریخ‌دار در
انتهای پرونده است که بگوید چه چیزی و چرا عوض شد، و هر نتیجه‌ای که زیر نسخهٔ پیشین محاسبه شده
برچسب خودش را نگه می‌دارد. این قاعده وجود دارد چون تغییر خاموش پروتکل برای بهتر کردن ظاهر یک
عدد، دقیقاً همان شکستی است که این پرونده برای جلوگیری از آن نوشته شده است.

\subsection{واحد استقلال}

واحد استقلال بیمار است، و در نبود شناسهٔ بیمار، چشم. دو چشم یک نفر و دو میدان یک چشم هرگز از
دو سوی یک تقسیم قرار نمی‌گیرند.

هیچ‌کدام از سه پیکرهٔ در دسترس شناسهٔ بیمار منتشر نمی‌کنند. پس پیش از نوشتن نخستین تقسیم، برای
هر تصویر یک درهم‌ساز ادراکی و یک توصیفگر الگوی عروقی محاسبه شد، خوشه‌ها ساخته شد، و هر خوشه یک
گروه در نظر گرفته شد. همین کار هم‌زمان تکراری‌های میان‌پیکره‌ای را پیدا کرد، که فرضی نیستند:
\num{۶۸۷} پرونده از \num{۱۷۴۴} پروندهٔ آینهٔ «\lr{Messidor-2}» ما نام‌گذاری به سبک \lr{Messidor-1} دارند.

تقسیم حاصل در \lr{data/splits/} نوشته و کامیت شد و هر مدل این پژوهش آن را از همان‌جا می‌خواند.
اثر انگشت تقسیم \lr{0cfbbfeb081999af} است و در نتایج هر اجرا ثبت می‌شود؛ اجرایی که اثر انگشت
متفاوتی داشته باشد با اجرای دیگر مقایسه نمی‌شود.

\subsection{آنچه کنار گذاشته شد}

دو سطح، چون یک تقسیم پانزده‌درصدی از \lr{IDRiD} چیزی را که دنبالش هستیم اندازه نمی‌گیرد.

\textbf{استخر توسعه} شامل \num{۵۱۶} تصویر \lr{IDRiD} و \num{۱۷۴۴} تصویر \lr{Messidor-2}، جمعاً \num{۲۲۶۰} تصویر
است و آموزش و انتخاب مدل هر دو روی آن انجام می‌شود، با اعتبارسنجی متقاطع گروهی پنج‌تایی.
پیکرهٔ \lr{EyePACS} تنها برای پیش‌آموزش به کار می‌رود و هرگز در انتخاب مدل دخالت نمی‌کند.

\textbf{مجموعهٔ آزمون بیرونی} \lr{APTOS 2019} \cite{aptos2019} است، با \num{۳۶۶۲} تصویر که به‌طور کامل کنار گذاشته
شده و تا انجماد مدل نهایی دیده نمی‌شود. دلیل انتخاب آن این است که از پیش در اختیار بود، پس
کنار گذاشتنش هزینه‌ای نداشت؛ به اندازهٔ کافی بزرگ است که بازهٔ اطمینان حدود یک واحد بدهد؛ و
جابه‌جایی توزیع واقعی است، نه بازبرش همان پیکره.

یک هشدار باید همراه هر عدد بیرونی بیاید: برچسب‌های \lr{APTOS} تک‌داور و نوفه‌دارترند. بخشی از
افت درون‌به‌بیرون نوفهٔ برچسب است، نه ناتوانی در تعمیم، و نباید همه‌اش به مدل نسبت داده شود.

\subsection{انتخاب مدل}

پیکربندی‌ها بر پایهٔ کارایی روی \emph{اعتبارسنجی} رتبه‌بندی می‌شوند، هرگز بر پایهٔ آزمون. این
شامل انتخاب \emph{میان} آزمایش‌ها هم می‌شود، نه فقط فراپارامترها درون یک آزمایش.

این قاعده در نرم‌افزار تثبیت شده است. اسکریپت تولید گزارش، مدل گزارش‌شده را از پروندهٔ
\lr{data/selected\_model.json} می‌خواند و اگر آن پرونده نباشد اجرا را رد می‌کند. پیش‌تر همان
اسکریپت مدل را با بیشینه‌گیری روی \lr{QWK} انتخاب می‌کرد، که انتخاب بر پایهٔ نمره است و وقتی
ارزیابی‌های بیرونی هم در استخر باشند، انتخاب روی مجموعهٔ آزمون می‌شود. اکنون اگر اجرای دیگری
نمرهٔ بالاتری داشته باشد، اسکریپت آن را با صدای بلند اعلام می‌کند و همچنان مدل انتخاب‌شده را
گزارش می‌کند.

\subsection{معیارها}

معیار اصلی برای رتینوپاتی دیابتی، \lr{QWK} است \cite{cohen1968kappa}. این معیار برای خطای دو درجه‌ای بیش از خطای یک
درجه‌ای جریمه می‌گیرد، که همان حالت شکستی است که هدف اصلاح آن را داریم. دقت همیشه در کنارش
گزارش می‌شود و هرگز به تنهایی. بازخوانی هر کلاس نیز همراه هر عدد سرتیتر می‌آید، چون مدل
می‌تواند دقت بگیرد و هم‌زمان نسبت به کلاس نادر کور شود.

برای تصمیم غربالگری، حساسیت و ویژگی رتینوپاتی نیازمند ارجاع گزارش می‌شود. این همان تصمیمی است
که سامانه از آن پشتیبانی می‌کند و همان چیزی است که بالینگر می‌پرسد.

برای ادم ماکولا، تعریف سه‌کلاسهٔ بدون دروازه اصلی است. جزئیات این تصمیم در بخش
\ref{sec:gating} می‌آید.

هر عدد بازهٔ خودگردان‌سازی دارد. چون سطرها همبسته‌اند، \textbf{بازنمونه‌گیری روی گروه‌ها انجام
می‌شود، نه روی تصویرها}. بازنمونه‌گیری روی تصویرها با نزدیک‌تکراری‌ها مانند شواهد مستقل رفتار
می‌کند و بازه را حدود دو برابر باریک‌تر از واقع نشان می‌دهد.

برای مقایسهٔ دو مدل، خودگردان‌سازی جفتی روی همان گروه‌ها انجام می‌شود و بازهٔ اختلاف گزارش
می‌شود. اگر بازه صفر را در بر بگیرد، دو مدل تفکیک‌ناپذیرند و همان‌طور گزارش می‌شوند؛ رتبه‌بندی
نمی‌شوند.

\subsection{واسنجی همتاشده پیش از انتساب}
\label{sec:matched}

\textbf{هیچ اختلافی میان دو مدل به بازنمایی نسبت داده نمی‌شود مگر آن‌که از آزمون واسنجی همتاشده
جان به در ببرد.} پیش از آن‌که بگوییم یک مدل ویژگی‌های بهتری آموخته، هر دو باید به یک شکل واسنجی
شوند — نقاط برش با برازش متقاطع و زیر یک هدف یکسان — و مقایسه تکرار شود.

دلیل وجود این قاعده تجربی است. آستانهٔ پیش‌فرض \num{۰٫۵} روی خروجی سیگموئید چیزی نیست که از مقیاس
درجه‌بندی آمده باشد؛ چیزی است که از تابع فعال‌سازی درمی‌آید. \textbf{یک پیش‌فرض تنظیم‌نشده،
فراپارامتری است که انتخاب شده، نه فراپارامتری که از آن پرهیز شده}، و هر مقایسه‌ای که آن را
میان مدل‌ها ثابت نگه دارد به آن آلوده است.

این قاعده تا امروز چهار ادعا را در \emph{دو جهت} برگردانده و دو بار علامت اختلاف را وارونه
کرده است. سابقه‌اش در جدول \ref{tab:matched_record} فصل بعد آمده است. نکتهٔ مهم این است که
قاعده‌ای که همیشه در یک جهت عمل کند آزمون نیست: سطری که نتیجهٔ مثبت داده همان چیزی است که به
سطرهای دیگر اعتبار می‌دهد.

\subsection{دو صلاحیت که همراه هر ادعای کلاسی می‌آیند}

نخست، \textbf{هیچ تصحیحی برای مقایسه‌های چندگانه اعمال نمی‌شود}. پنج کلاس ضرب در شرایط
مقایسه‌شده یعنی به‌طور متوسط یک خانهٔ معنادار از سر تصادف انتظار می‌رود. پس اعتبار از
\emph{تکرار یک اختلاف در شرایط مختلف} می‌آید، نه از مقدار \lr{p} به تنهایی.

دوم، جایی که مقایسه بیش از یک چیز را عوض می‌کند، ادعای کلاسی همهٔ مخدوش‌کننده‌های ادعای کلی را
به ارث می‌برد.

\section{داده}

\subsection{پیکره‌ها}

\lr{IDRiD} \cite{porwal2018idrid} با \num{۵۱۶} تصویر تنها منبع سه‌کلاسهٔ ادم ماکولا در این پژوهش است و
\textbf{\num{۵۱} تصویر در درجهٔ میانی} دارد. همین عدد سقف کارایی سر ادم ماکولا را تعیین می‌کند و در
فصل \ref{ch:results} به آن بازمی‌گردیم.

\lr{Messidor-2} \cite{decenciere2014messidor} با \num{۱۷۴۴} تصویر برچسب رتینوپاتی دارد و برای ادم ماکولا تنها برچسب دودویی
«نیازمند ارجاع یا نه» می‌دهد. این برچسب جزئی است، نه ناقص: در صورت‌بندی آستانه‌ای، «نیازمند
ارجاع» دقیقاً معادل درجهٔ دوی \lr{IDRiD} است، پس آستانهٔ دوم را دقیق نظارت می‌کند و آستانهٔ
نخست را پوشانده می‌گذارد. هیچ تقریبی در کار نیست.

\lr{APTOS 2019} \cite{aptos2019} با \num{۳۶۶۲} تصویر به‌طور کامل کنار گذاشته شده است.

\subsection{ارزیابی بدون دروازه به‌عنوان تعریف اصلی}
\label{sec:gating}

کار پیشین شاخهٔ ادم ماکولا را مشروط به رتینوپاتی بزرگ‌تر یا مساوی یک می‌کرد. در این پژوهش
ارزیابی \textbf{بدون دروازه} اصلی است و عدد مشروط به‌عنوان ثانویه و همیشه برچسب‌خورده گزارش
می‌شود. سه دلیل دارد.

نخست، \textbf{دروازه چیزی را کنار نمی‌گذارد}. روی \lr{IDRiD} بررسی شد: از \num{۵۱۶} تصویر، هر \num{۱۶۸}
تصویری که رتینوپاتی صفر دارند، درجهٔ ادم ماکولایشان هم صفر است. هیچ نمونهٔ مثبت ادم ماکولا با
رتینوپاتی صفر وجود ندارد. پس دروازه از چیزی محافظت نمی‌کند و تنها مخرج کسر را عوض می‌کند.

دوم، معیار مشروط به درست بودن سر رتینوپاتی وابسته است، پس میان مدل‌هایی با سرهای متفاوت
قابل مقایسه نیست.

سوم، و مهم‌تر از همه، \textbf{کف کلاس اکثریت به‌شدت فرق می‌کند}: \num{۶۹٫۸} درصد برای مجموعهٔ مشروط
در برابر \num{۴۷٫۱} درصد برای بدون دروازه. \textbf{هر عدد ادم ماکولا در این پژوهش کنار کف خودش نقل
می‌شود.}

\section{معماری و آموزش}

ستون فقرات مشترک است و دو سر جدا دارد، یکی برای رتینوپاتی با پنج درجه و یکی برای ادم ماکولا با
سه درجه. سرها آستانه‌ای‌اند: به جای بیشینهٔ نرم روی \lr{K} کلاس، \lr{K-1} احتمال از شکل
$P(y>k)$ تخمین زده می‌شود. این انتخاب سه کار می‌کند: ترتیب درجه‌ها را وارد تابع زیان می‌کند
که همان چیزی است که \lr{QWK} می‌سنجد؛ رمزگشایی را با ساخت سازگار با رتبه می‌کند؛ و از همه
مهم‌تر، برچسب جزئی \lr{Messidor-2} را قابل استفاده می‌کند.

درجه‌بندی پنج‌درجه‌ای رتینوپاتی و سه‌درجه‌ای ادم ماکولا از مقیاس بالینی بین‌المللی
می‌آید \cite{wilkinson2003scale}. ستون فقرات پیش‌فرض \lr{DenseNet121}
\cite{huang2017densenet} است و مدل انتخاب‌شده \lr{EfficientNet-B3}
\cite{tan2019efficientnet}؛ \lr{ConvNeXt-tiny} \cite{liu2022convnext} نیز به‌عنوان یک نامزد
آزموده شد. دو سر ترتیبی جایگزین که در فصل \ref{ch:results} ابطال می‌شوند
\lr{CORAL} \cite{cao2020rank} و \lr{CORN} \cite{shi2023corn} هستند، و راهبرد
\lr{LP-FT} از \cite{kumar2022lpft} می‌آید. گزارش این پژوهش در برابر سیاههٔ
\lr{CLAIM} \cite{mongan2020claim} و \lr{TRIPOD} \cite{collins2015tripod} سنجیده شده است، و
معیار واسنجی و مسئلهٔ بیش‌اطمینانی شبکه‌های عمیق از \cite{guo2017calibration} گرفته شده است.

پیش‌آموزش روی \lr{EyePACS} \cite{eyepacs2015} انجام می‌شود، یک بار، و همهٔ تاها از همان وزن
آغاز می‌کنند.

افزایش داده شامل چرخش آزاد، بازتاب، و لرزش ملایم رنگ و گاما است. چرخش نامحدود است چون شبکیه
جهت متعارف ندارد و هر معیار درجه‌بندی نسبت به چرخش ناوردا است.

\section{تبار داده}
\label{sec:provenance}

این بخش سهمی است، نه اعتراف. شش اشکال تبار داده در این پژوهش هیچ خطای زمان اجرا تولید نکردند و
در سکوت عدد غلط ساختند. سازوکارهای زیر پاسخ آن‌هاست.

\subsection{پیکربندی، مصرف نیست}

چهار سازوکار مستقل تبار داده هم‌زمان از اجرایی عبور کردند که ترکیبش غلط می‌بود: بلوک پیکربندی
یکسان، اثر انگشت تقسیم یکسان، نگهبان اسکریپت، و میدانی که برای هر دو اجرا \emph{به‌راستی درست}
بود و برای هرکدام معنای دیگری داشت. هر چهار سازوکار ثبت می‌کنند که اجرا چه چیزی را
\emph{پیکربندی} کرده است؛ هیچ‌کدام ثبت نمی‌کنند چه چیزی را \emph{مصرف} کرده است.

پس هر اجرا اکنون یک \textbf{بیانیهٔ مصرف} می‌نویسد: برای هر پیکره، تعداد تصویرهای بازیابی‌شده و
یک درهم‌ساز روی مسیرهای بازیابی‌شده نسبت به محل سوارشدن دیتاست. دو اجرا تنها وقتی ترکیب،
هم‌گروه یا مقایسه می‌شوند که بیانیه‌هایشان بخواند. مقایسه‌های عمداً میان‌منبعی ممکن‌اند، اما
پرچم لازم اختلاف را \emph{در سند خروجی مهر می‌کند}؛ هشداری که خواننده نبیند هشدار نیست.

\subsection{پنج نمونه، که یکی از آن‌ها خودِ ماست}
\label{sec:provenance-exhibits}

استدلال این بخش وقتی وزن دارد که نمونه داشته باشد. چهار نمونه از کارهای منتشرشده می‌آید و
\textbf{نمونهٔ پنجم از همین پایان‌نامه است}. ترتیب عمدی است: فصلی که تبار دادهٔ دیگران را
حسابرسی کند و شکست خودش را از همان جنس گزارش نکند، ادعایی می‌کند که خودش را زیر آن نمی‌برد.

\textbf{نمونهٔ یکم — سطح بالای \lr{AUC} به‌جای واسنجی \cite{yu2026dualswinord}.} یک مقالهٔ
داوری‌شدهٔ \num{۲۰۲۶} می‌نویسد که مقادیر بالای \lr{AUC} «تأیید می‌کنند» مدل احتمال‌های
خوش‌واسنجی می‌دهد. آن مقاله \emph{هیچ} معیار واسنجی گزارش نمی‌کند: نه \lr{ECE}، نه
\lr{Brier}، نه منحنی اطمینان. و \lr{AUC} نسبت به هر تبدیل یکنوای نمره‌ها ناوردا است، پس
دربارهٔ واسنجی هیچ نمی‌گوید؛ مدلی می‌تواند همان \lr{AUC} را نگه دارد و به هر اندازه
بدواسنجی باشد.

\textbf{نمونهٔ دوم — تحلیل چشمی روی داده‌های جفت‌شده \cite{tomic2024handheld}.} پژوهشی با \num{۳۲۰}
چشم از \num{۱۶۰} بیمار، همهٔ تحلیل‌ها را چشم‌به‌چشم انجام می‌دهد و برای همبستگی دو چشم یک نفر تصحیحی
نمی‌گذارد. بازه‌های اطمینان گزارش‌شده از این راه باریک‌تر از واقع درمی‌آیند. این همان دلیلی است
که در بخش \ref{sec:matched} بازنمونه‌گیری این پژوهش روی گروه‌ها انجام می‌شود، نه روی تصویرها.

\textbf{نمونهٔ سوم — ستونی که از ستون دیگر رونویسی شده \cite{tahir2025meta}.} در یک
فراتحلیل، ستون منفی‌های کاذب عیناً تکرار ستون مثبت‌های کاذب است. این خطا در جدول نمی‌ماند: از
همان اعداد، حساسیت یکی از پژوهش‌های واردشده \num{۰٫۰۶} درمی‌آید، در حالی که حساسیت تجمیعی گزارش‌شده
\num{۰٫۹۵} است. \textbf{خرابی داده به عددهای گزارش‌شده نشت کرده است.}

\textbf{نمونهٔ چهارم — تقسیمی که در دو پیمانه اعمال می‌شود و در سومی نه
\cite{tang2026mmrdr}.} یک توصیف‌نامهٔ داده در خانوادهٔ \lr{Nature} زیرمجموعه‌های خود را در دو
پیمانه در سطح بیمار تقسیم می‌کند «برای جلوگیری از نشت داده»، و زیرمجموعهٔ سوم را در سطح تصویر،
چون پیکرهٔ مبدأ شناسهٔ بیمار ندارد؛ سپس یک جدول محک واحد روی هر سه گزارش می‌شود. قاعده
شناخته شده، نوشته شده، و در یک‌سوم موارد اجرا نشده است.

\subsection{نمونهٔ پنجم: شکست خودِ این متن}
\label{sec:our-own-failure}

فصل حاضر در نسخهٔ نخست خود با این جمله آغاز می‌شد که همهٔ عددهایش از اجراهای آرشیوشده می‌آیند و
هیچ عددی با دست تایپ نشده است. \textbf{این جمله درست نبود.} عددها با دست تایپ شده بودند، هیچ
سازوکاری آن‌ها را بررسی نمی‌کرد، و سه عدد غلط بود.

\begin{table}[H]
\caption{سه خطای واقعی در نسخه‌های نخست همین متن، که تنها به این دلیل پیدا شدند که ساختن یک
دفتر اعداد، حل کردن هر رقم در برابر یک مصنوع آرشیوی را ضروری کرد.}
\centering
\begin{tabular}{|l|c|c|}
\hline
\textbf{عدد} & \textbf{نوشته‌شده} & \textbf{درست} \\ \hline
کف مشروط ادم ماکولا & \num{۶۹٫۶}٪ & \num{۶۹٫۸}٪ \\ \hline
تطابق دقیق در فصل \ref{sec:parta} & \num{۹۶٫۶}٪ & \num{۹۶٫۳}٪ \\ \hline
اختلاف هم‌نهاد تصویرها & \num{۰٫۷۸} تا \num{۰٫۹۶} & \num{۰٫۶۵} تا \num{۰٫۷۵} \\ \hline
\end{tabular}
\label{tab:our-failure}
\end{table}

\textbf{سطر سوم یافته است، نه غلط تایپی.} آن بازه شاهدِ نخستین توضیح رقیبی بود که فصل
\ref{sec:parta} باید رد می‌کرد: این‌که تناظر میان شماره‌گذاری قطعه‌بندی و شماره‌گذاری
درجه‌بندی، ماسک را کنار درجهٔ اشتباه گذاشته باشد. تمام آن حکم روی این عدد می‌ایستاد.

\textbf{و آن عدد هیچ مصنوعی نداشت.} نه اسکریپتی، نه پروندهٔ خروجی، نه ثبتی از روش. بازتولید
نشد، و وقتی از نو محاسبه شد پاسخ \num{۰٫۶۵} تا \num{۰٫۷۵} درآمد. نتیجه عوض نشد — جفت‌ها همان
عکس‌اند و حکم سرجایش است — ولی \textbf{عددهای منتشرشده بازتولیدپذیر نبودند و با عددهایی
جایگزین شدند که هستند.} عددی که اسکریپتش دیگر وجود ندارد، ادعا است نه اندازه‌گیری؛ و این همان
جمله‌ای است که این فصل دربارهٔ مقاله‌های دیگران می‌گوید.

\subsection{و نمونهٔ ششم، که هنگام نوشتن نمونهٔ پنجم ساخته شد}

سه خطای بالا هنگام بررسی گسترده‌تری پیدا شدند، و آن بررسی خودش به یک اشتباه انجامید که آموزندهٔ
هر پنج نمونهٔ پیشین است.

مشاهده این بود که عددهای اعشاری فصل‌های نخست، بخش صحیح و بخش اعشاری‌شان جابه‌جا نوشته شده است.
تشخیص از روی استاندارد ساخته شد: ارقام فارسی در جدول یونیکد از ردهٔ عدد اروپایی‌اند، رشته‌های
عددی چپ‌به‌راست چیده می‌شوند، پس متن همان‌طور که نوشته شده چاپ می‌شود، پس
\textbf{فصل روش به خواننده می‌گفت آستانهٔ پیش‌فرض \wrongnum{۵/۰} است}.

\textbf{استدلال دربارهٔ یونیکد درست بود و نتیجه دربارهٔ این سند غلط.} پیاده‌سازی دوسویهٔ
موتور حروف‌چینی این پایان‌نامه، پیاده‌سازی کامل الگوریتم یونیکد نیست. عددی که به ترتیب منطقی
نوشته شود، \emph{وارونه} چاپ می‌شود. \textbf{پس نسخه‌های نخست درست چاپ می‌شدند}؛ نویسنده‌شان
برای موتور چاپ می‌نوشت. «تصحیح» منبع را منطقی کرد و \textbf{صفحهٔ چاپ‌شده را غلط}، و همین‌طور
ماند تا وقتی که سرانجام یک صفحه چاپ و نگاه شد.

\textbf{جنس این خطا دقیقاً همان جنس چهار نمونهٔ منتشرشده است.} در نمونهٔ نخست، ادعای واسنجی
جای اندازه‌گیری واسنجی را گرفته بود. اینجا، مشخصات یک استاندارد جای نگاه کردن به خروجی را
گرفت. \textbf{هیچ‌کدام به خودِ چیز نگاه نکردند.}

قرارداد عوض شد، ولی \emph{به دلیل خودش}، نه به‌عنوان رفع اشکال: منبع به ترتیب منطقی نوشته
می‌شود و هر عدد در یک بازهٔ صریح چپ‌به‌راست پیچیده می‌شود. این هم درست چاپ می‌شود و هم
جست‌وجوپذیر و بررسی‌پذیر است؛ قرارداد پیشین درست چاپ می‌شد و هیچ‌کدام از این دو نبود.

\subsection{آنچه اکنون این را نگه می‌دارد}

هر عدد فارسیِ متن باید در \lr{docs/generated/thesis\_numbers.json} باشد، که از مصنوع‌های
آرشیوی ساخته می‌شود. ویراستار خودکار سه چیز را رد می‌کند: رقمی که در آن دفتر نباشد، جداکنندهٔ
اعشار نادرست، و \textbf{عددی که در پوشش چپ‌به‌راست نباشد}. آزمون سوم دقیقاً به این دلیل وجود
دارد که نمونهٔ ششم از دو آزمون نخست عبور کرد.

هر سه آزمون پیش از اعتماد، در هر دو جهت عمل کرده‌اند. \textbf{آنچه ثابت نمی‌کنند} این است که
عدد برای \emph{آن جمله} درست باشد؛ هیچ آزمون خودکاری آن را ثابت نمی‌کند، و ادعای خلافش دقیقاً
همان «آزمونی که موفقیت گزارش می‌کند بی‌آن‌که چیزی را بررسی کرده باشد» است.

\textbf{درسی که می‌ماند ساده است:} ادعایی دربارهٔ خروجی، با نگاه کردن به خروجی حل می‌شود. دو
نگهبان این پژوهش به این دلیل وجود دارند که یک بار استدلال جای بازرسی را گرفت.

\subsection{آزمونی که فقط در یک جهت عمل کند آزمون نیست}

هر قاعده در این پژوهش سابقه‌اش را حمل می‌کند، و سابقه باید ابطال‌هایی در هر دو جهت داشته باشد.
قاعده‌ای که تنها نتایج ناخوشایند را حل کرده باشد از «راهی برای توجیه» قابل تشخیص نیست، و
قاعده‌ای که تنها تأیید کرده باشد مهر لاستیکی است.

حالت انحطاطی این وضع بدترین شکست ممکن برای یک نگهبان است: \textbf{آزمونی که موفقیت گزارش کند
بی‌آن‌که چیزی را بررسی کرده باشد.} دو نمونه از آن در این پژوهش رخ داد و هر دو اکنون ادعا
می‌کنند که چیزی بازرسی شده است. نگهبانی که با صدای بلند شکست بخورد یک اجرا هزینه دارد؛
نگهبانی که به غلط قبول شود، اعتماد به این‌که اصلاً بررسی کرده‌اید را هزینه دارد.

% فصل ۴ — نتایج
% همهٔ عددها از docs/generated/thesis_numbers.json می‌آیند، که src/thesis_numbers.py آن را
% از آرشیو runs/ و docs/generated/ می‌سازد. صحت با «python3 tools/farsi_lint.py --numbers» بررسی می‌شود.

\chapter{نتایج}
\label{ch:results}

این فصل با پرسشی آغاز می‌شود که معمولاً پرسیده نمی‌شود: خط لولهٔ آموزش در برابر \emph{چه چیزی}
سنجیده می‌شود؟ سپس نتیجهٔ مدل انتخاب‌شده روی استخر توسعه و مجموعهٔ بیرونی می‌آید، بعد سر ادم
ماکولا و سقفی که به آن می‌خورد، و در پایان فهرست چیزهایی که کار نکردند. این ترتیب عمدی است:
هشت مداخله از نه مداخلهٔ آزموده‌شده ابطال شدند، و پنهان کردن آن‌ها در پیوست، تصویری از پژوهش
می‌سازد که با آنچه رخ داد نمی‌خواند.

\section{خط لوله در برابر چه چیزی سنجیده می‌شود}
\label{sec:baseline}

پیش از هر مقایسه‌ای میان مدل‌های این پژوهش، یک عدد لازم است: خط لولهٔ آموزش چقدر بر
ساده‌ترین جایگزین ممکن می‌افزاید. بدون آن، هر مقایسه‌ای میان دو مدلِ \emph{هر دو خوب} انجام
می‌شود و این پیش‌فرض را می‌پذیرد که ساختن خط لوله ارزشش را داشته است.

سنجه روی همان تصویرها، همان تاها، همان رمزگشایی و همان واسنجی همتاشده انجام شد. پایه، ویژگی‌های
منجمد \lr{DenseNet121} از پیش‌آموزش \lr{ImageNet} است با یک سر خطی ترتیبی روی آن.

\begin{table}[H]
\caption{ارزش خط لولهٔ کامل در برابر ویژگی‌های منجمد و یک مدل خطی. هر دو سطر روی \num{۲۲۶۰} تصویر
یکسان، تاهای یکسان و واسنجی همتاشده محاسبه شده‌اند.}
\centering
\begin{tabular}{|l|c|}
\hline
\textbf{مدل} & \textbf{\lr{QWK}} \\ \hline
پیش‌بین کلاس اکثریت & \num{۰} به‌طور ساختاری \\ \hline
ویژگی‌های منجمد \lr{ImageNet} + سر خطی ترتیبی & \num{۰٫۶۶۱۱} \\ \hline
ویژگی‌های منجمد \lr{RETFound} + سر خطی ترتیبی & \num{۰٫۷۰۰۸} \\ \hline
\textbf{خط لولهٔ کامل (مدل انتخاب‌شده)} & \textbf{\num{۰٫۸۹۵۴}} \\ \hline
\end{tabular}
\label{tab:baseline}
\end{table}

\textbf{خط لولهٔ کامل \num{۰٫۲۳۴۲} واحد \lr{QWK} بر ویژگی‌های منجمد و یک مدل خطی می‌افزاید.} این یک
گزارهٔ همتا به همتا است و نخستین عددی است که در این پژوهش به پرسش «در مقایسه با چه؟» بدون خروج
از تقسیم منجمد پاسخ می‌دهد.

دو صلاحیت همراه آن می‌آید. نخست، این مقایسه با کار منتشرشده نیست و نمی‌تواند بشود؛ تقسیم
اختصاصی و منجمد هر عدد منتشرشده‌ای را ناهم‌سنجه می‌کند. دوم، ادعا این نیست که \num{۰٫۶۶۱۱} پایهٔ
\emph{خوبی} است — ویژگی‌های منجمد \lr{ImageNet} روی تصویر شبکیه سطح پایینی‌اند. نکته فاصله از
آن است، که اندازه گرفته شد و فرض نشد.

\section{مدل انتخاب‌شده}
\label{sec:champion}

مدل گزارش‌شده بر پایهٔ رتبه‌بندی روی استخر توسعه انتخاب شد و در
\lr{data/selected\_model.json} اعلام شده است. این یک واقعیت اعلام‌شده است، نه استنتاجی از
نمره‌ها: اسکریپت گزارش مدل را از آن پرونده می‌خواند و اگر اجرای دیگری نمرهٔ بالاتری بگیرد، آن
را با صدای بلند اعلام می‌کند و همچنان مدل اعلام‌شده را گزارش می‌کند.

\subsection{رتینوپاتی دیابتی روی استخر توسعه}

\begin{table}[H]
\caption{کارایی به تفکیک کلاس برای درجه‌بندی پنج‌کلاسهٔ رتینوپاتی دیابتی، برون‌تایی روی
\num{۲۲۶۰} تصویر استخر توسعه، با نقاط برشِ برازش‌متقاطع.}
\centering
\begin{tabular}{|l|c|c|}
\hline
\textbf{درجه} & \textbf{تعداد} & \textbf{بازخوانی} \\ \hline
\num{۰} — بدون رتینوپاتی & \num{۱۱۸۵} & \num{۹۴٫۲}٪ \\ \hline
\num{۱} — خفیف & \num{۲۹۵} & \textbf{\num{۳۱٫۹}٪} \\ \hline
\num{۲} — متوسط & \num{۵۱۵} & \num{۷۱٫۱}٪ \\ \hline
\num{۳} — شدید & \num{۱۶۸} & \num{۷۵٫۰}٪ \\ \hline
\num{۴} — تکثیری & \num{۹۷} & \num{۵۸٫۸}٪ \\ \hline
\textbf{کل} & \num{۲۲۶۰} & \multicolumn{1}{c|}{دقت \num{۷۷٫۸}٪، \lr{QWK} \num{۰٫۸۹۵۴}} \\ \hline
\end{tabular}
\label{tab:dr-dev}
\end{table}

عدد سرتیتر \num{۰٫۸۹۵۴} است و ستون بازخوانی چیزی می‌گوید که عدد سرتیتر نمی‌گوید: درجهٔ خفیف با
بازخوانی \num{۳۱٫۹} درصد بازیابی می‌شود، یعنی حدود دو سوم موارد خفیف به درجهٔ همسایه می‌روند. این
همان دلیلی است که در فصل \ref{ch:method} گفته شد بازخوانی هر کلاس همراه هر عدد سرتیتر می‌آید:
یک مدل می‌تواند دقت خوبی بگیرد و هم‌زمان نسبت به یک کلاس تقریباً کور بماند.

\textbf{اما این کوری، ظرفیت نیست؛ واسنجی است.} بخش \ref{sec:matched-record} نشان می‌دهد که
جابه‌جایی نقاط برش، بدون آموزش دوباره و بدون تغییر مدل، همین بازخوانی را چند برابر می‌کند. این
مهم‌ترین یافتهٔ روشی این فصل است و پیش از هر تفسیر معماری باید خوانده شود.

\subsection{اعتبارسنجی بیرونی}

مجموعهٔ \lr{APTOS} تا انجماد مدل نهایی دیده نشد و در هیچ مرحله‌ای از آموزش یا انتخاب مدل
دخالت نداشت. پیش‌بینی‌ها میانگین لاجیت پنج تا با افزایش داده در زمان آزمون است.

\begin{table}[H]
\caption{اعتبارسنجی بیرونی مدل انتخاب‌شده روی \num{۳۶۶۲} تصویر \lr{APTOS}، که در هیچ مرحله‌ای از
آموزش یا انتخاب مدل دیده نشده است.}
\centering
\begin{tabular}{|l|c|}
\hline
\textbf{معیار} & \textbf{مقدار} \\ \hline
تعداد تصویرها & \num{۳۶۶۲} \\ \hline
کف کلاس اکثریت & \num{۴۹٫۳}٪ \\ \hline
دقت & \num{۷۴٫۹}٪ \quad [\num{۷۳٫۵}، \num{۷۶٫۳}] \\ \hline
\lr{QWK} & \textbf{\num{۰٫۹۰۲۶}} \quad [\num{۰٫۸۹۴۱}، \num{۰٫۹۱۰۸}] \\ \hline
میانگین کلان \lr{F1} & \num{۰٫۵۴۵} \\ \hline
حساسیت رتینوپاتی نیازمند ارجاع & \num{۹۹٫۳}٪ \\ \hline
ویژگی رتینوپاتی نیازمند ارجاع & \num{۸۴٫۳}٪ \\ \hline
\end{tabular}
\label{tab:external}
\end{table}

\lr{QWK} بیرونی \num{۰٫۹۰۲۶} است، اندکی \emph{بالاتر} از عدد درونی \num{۰٫۸۹۵۴}. این را نباید به‌عنوان
تعمیم بهتر خواند. سه چیز هم‌زمان فرق می‌کنند: توزیع درجه‌ها، شیوهٔ برچسب‌گذاری، و اندازهٔ
نمونه. \lr{APTOS} برچسب تک‌داور دارد و نوفهٔ برچسبش بیشتر است؛ بخشی از هر تفاوت درون‌به‌بیرون
نوفهٔ برچسب است، نه توانایی مدل.

\textbf{و ستون بازخوانی، همان الگو را با شدت بیشتر تکرار می‌کند.}

\begin{table}[H]
\caption{بازخوانی به تفکیک کلاس روی \lr{APTOS}. کل \lr{QWK} بالاست و درجهٔ خفیف تقریباً
به‌طور کامل از دست می‌رود.}
\centering
\begin{tabular}{|l|c|c|}
\hline
\textbf{درجه} & \textbf{تعداد} & \textbf{بازخوانی} \\ \hline
\num{۰} — بدون رتینوپاتی & \num{۱۸۰۵} & \num{۹۸٫۱}٪ \\ \hline
\num{۱} — خفیف & \num{۳۷۰} & \textbf{\num{۴٫۹}٪} \\ \hline
\num{۲} — متوسط & \num{۹۹۹} & \num{۶۵٫۶}٪ \\ \hline
\num{۳} — شدید & \num{۱۹۳} & \num{۶۳٫۲}٪ \\ \hline
\num{۴} — تکثیری & \num{۲۹۵} & \num{۶۰٫۳}٪ \\ \hline
\end{tabular}
\label{tab:external-recall}
\end{table}

\lr{QWK} برابر \num{۰٫۹۰۲۶} در کنار بازخوانی \num{۴٫۹} درصد برای درجهٔ خفیف می‌نشیند. \textbf{این تناقض
نیست؛ این همان چیزی است که \lr{QWK} می‌سنجد.} خطای یک درجه‌ای کم جریمه می‌شود، و درجهٔ خفیف
دقیقاً میان دو کلاس پرجمعیت‌تر قرار دارد، پس فروریختن آن به همسایه‌ها به عدد سرتیتر آسیب کمی
می‌زند. هر گزارشی که تنها \lr{QWK} را نقل کند، این را نشان نمی‌دهد.

از منظر تصمیم غربالگری، عددی که اهمیت دارد حساسیت رتینوپاتی نیازمند ارجاع است: \num{۹۹٫۳} درصد با
ویژگی \num{۸۴٫۳} درصد. درجهٔ خفیف \emph{نیازمند ارجاع نیست}، پس فروریختن آن به درجهٔ صفر تصمیم ارجاع
را خراب نمی‌کند — اما درجه‌بندی پنج‌کلاسه را خراب می‌کند، و پایان‌نامه هر دو را ادعا می‌کند.

\section{ادم ماکولای دیابتی و سقفی که به آن می‌خورد}
\label{sec:dme}

هر عدد ادم ماکولا کنار کف کلاس اکثریت خودش نقل می‌شود. ارزیابی بدون دروازه اصلی است و
دلیلش در بخش \ref{sec:gating} آمد.

\begin{table}[H]
\caption{کارایی به تفکیک کلاس برای درجه‌بندی سه‌کلاسهٔ ادم ماکولا، بدون دروازه، برون‌تایی روی
\num{۵۱۶} تصویر \lr{IDRiD} — تنها منبع سه‌کلاسهٔ در دسترس.}
\centering
\begin{tabular}{|l|c|c|}
\hline
\textbf{درجه} & \textbf{تعداد} & \textbf{بازخوانی} \\ \hline
\num{۰} — بدون ادم ماکولا & \num{۲۲۲} & \num{۸۸٫۳}٪ \\ \hline
\num{۱} — غیرقابل‌ارجاع & \num{۵۱} & \textbf{\num{۵۴٫۹}٪} \\ \hline
\num{۲} — قابل‌ارجاع & \num{۲۴۳} & \num{۹۰٫۵}٪ \\ \hline
\textbf{کل} & \num{۵۱۶} & \multicolumn{1}{c|}{دقت \num{۸۶٫۰}٪، \lr{QWK} \num{۰٫۸۸۸۲}} \\ \hline
\textit{کف کلاس اکثریت} & & \multicolumn{1}{c|}{\textit{\num{۴۷٫۱}٪}} \\ \hline
\end{tabular}
\label{tab:dme-dev}
\end{table}

الگو همان است: دو کلاس بیرونی خوب بازیابی می‌شوند و کلاس میانی نصف می‌شود.
\textbf{و اینجا علتش را می‌دانیم.} تنها \num{۵۱} تصویر در کل پژوهش درجهٔ میانی ادم ماکولا دارند.
این عدد از معماری نمی‌آید؛ از داده می‌آید.

سه مداخلهٔ مستقل این را آزمودند و هر سه ابطال شدند: برش مرکز-ماکولا، رزولوشن بومی، و ترکیب
مدل‌ها. هیچ‌کدام بازهٔ اطمینان را از صفر دور نکردند.

برش مرکز-ماکولا از این میان مهم‌ترین است، چون دقیقاً همان ناحیه‌ای را جدا می‌کند که تعریف
بالینی درجه بر آن بنا شده است. \textbf{وقتی مداخله‌ای که ناحیهٔ تعریف‌کنندهٔ برچسب را هدف
می‌گیرد کاری نمی‌کند، محدودیت در نظارت است، نه در بازنمایی.}

پیامد عملی این است که سر ادم ماکولا با داده‌ای که در دسترس است بهتر نمی‌شود، و فصل
\ref{sec:parta} نشان می‌دهد چرا داده‌ای هم به آن افزوده نشد.

\section{آنچه کار نکرد}
\label{sec:falsified}

هر مداخله پیش از اجرا فرضیه و پیامد ابطال‌گر خود را ثبت کرد. جدول \ref{tab:falsified} همهٔ
آن‌ها را می‌آورد، با \lr{QWK} رتینوپاتی زیر واسنجی همتاشده و بازهٔ جفتی خودگردان‌سازی روی
گروه‌ها.

\begin{table}[H]
\caption{نه مداخله، یک تأیید. اختلاف‌ها زیر واسنجی همتاشده و با خودگردان‌سازی جفتی روی
گروه‌ها. بازه‌ای که صفر را در بر بگیرد یعنی دو حالت تفکیک‌ناپذیرند. ستون «پایه» می‌گوید هر
سطر روی چند تصویر برون‌تایی سنجیده شده است، چون سه سطر آخر روی همهٔ استخر اجرا نشدند.}
\centering
\begin{tabular}{|l|c|c|c|c|}
\hline
\textbf{مداخله} & \textbf{اختلاف \lr{QWK}} & \textbf{بازهٔ \num{۹۵} درصد} & \textbf{پایه} & \textbf{حکم} \\ \hline
\lr{EfficientNet-B3} به جای \lr{DenseNet121} & \num{۰٫۰۲۰۷} & [\num{۰٫۰۱۰۵}، \num{۰٫۰۳۲۰}] & \num{۲۲۶۰} & \textbf{تأیید} \\ \hline
رزولوشن بومی & \num{۰٫۰۰۸۶} & [-\num{۰٫۰۰۲۵}، \num{۰٫۰۱۹۹}] & \num{۲۲۶۰} & ابطال \\ \hline
برش مرکز-ماکولا & \num{۰٫۰۰۴۹} & [-\num{۰٫۰۰۷۸}، \num{۰٫۰۱۷۱}] & \num{۲۲۶۰} & ابطال \\ \hline
\lr{LP-FT} & -\num{۰٫۰۰۰۸} & [-\num{۰٫۰۱۲۶}، \num{۰٫۰۱۰۵}] & \num{۲۲۶۰} & ابطال \\ \hline
سر ترتیبی \lr{CORAL} & -\num{۰٫۰۰۷۶} & [-\num{۰٫۰۱۸۴}، \num{۰٫۰۰۳۲}] & \num{۲۲۶۰} & ابطال \\ \hline
سر ترتیبی \lr{CORN} & -\num{۰٫۰۴۷۹} & [-\num{۰٫۰۶۵۶}، -\num{۰٫۰۳۲۵}] & \num{۲۲۶۰} & ابطال \\ \hline
ترکیب مدل‌ها (روی مجموعهٔ بیرونی) & \num{۰٫۰۰۲۵} & [-\num{۰٫۰۰۱۶}، \num{۰٫۰۰۶۶}] & \num{۳۶۶۲} & ابطال \\ \hline
\lr{ConvNeXt-tiny} به جای \lr{DenseNet121} & -\num{۰٫۰۳۶۲} & [-\num{۰٫۰۵۹۳}، -\num{۰٫۰۱۵۳}] & \num{۹۱۰} & ابطال \\ \hline
ریزتنظیم مدل بنیادین شبکیه & -\num{۰٫۰۳۹۹} & [-\num{۰٫۰۶۶۹}، -\num{۰٫۰۱۲۴}] & \num{۹۱۰} & ابطال \\ \hline
\end{tabular}
\label{tab:falsified}
\end{table}

سه سطر با اختلاف معنادار در جهت \emph{نامطلوب} پایان می‌یابند. پنج سطر بازه‌شان صفر را در بر
می‌گیرد و تفکیک‌ناپذیر گزارش می‌شوند، نه «کمی بهتر». تنها یک مداخله از نه مداخله تأیید شد.

\textbf{سه سطر آخر روی همهٔ استخر اجرا نشدند و به همان صورت خوانده می‌شوند.} ترکیب مدل‌ها روی
مجموعهٔ بیرونی سنجیده شد، چون همان‌جا بود که ادعایش شکست: سودی که روی استخر توسعه داشت به
مجموعهٔ کنارگذاشته منتقل نشد. دو سطر دیگر روی دو تا از پنج تا اجرا شدند و وقتی نتیجه در \num{۴۰}
درصد داده معنادار و منفی بود متوقف شدند — صرف ساعت‌های \lr{GPU} برای تیزتر کردن عددی که پیشاپیش
منفی است، تصمیم بدی است، و متوقف کردنشان با همان قاعده‌ای انجام شد که پیش از اجرا ثبت شده بود.

\textbf{سطر آخر با \lr{--acknowledge-consumption-diff} اجرا شد، و این عمدی است.} مدل بنیادین
جای مرحلهٔ پیش‌آموزش \lr{EyePACS} را می‌گیرد، پس دو اجرا به‌طور ساختاری پرونده‌های متفاوتی
می‌خوانند. این تفاوت خودِ آزمایش است، ولی پرچمِ لازم آن را در سند خروجی مهر می‌کند؛ هشداری که
خواننده نبیند هشدار نیست.

دو نکتهٔ روشی همراه این جدول می‌آید.

\textbf{نخست، خطِ زیان ترتیبی با دو اجرا بسته شد، نه با یکی.} نتیجهٔ منفی \lr{CORAL} به تنهایی
مبهم بود، چون \lr{CORAL} هم‌زمان دو چیز را عوض می‌کند: هدف آموزش، و محدود کردن سر به یک جهت
مشترک به جای \lr{K-1} تصویرسازی مستقل. پس \lr{CORN} اجرا شد، که همان چارچوب ترتیبی را بدون آن
محدودیت ظرفیت نگه می‌دارد. نتیجه‌اش به‌روشنی بدتر بود، و علتش قابل اندازه‌گیری است:
زیرمجموعه‌های شرطی \lr{CORN} زیر نامتوازنی کلاس‌ها فرومی‌ریزند و آستانهٔ سوم تنها روی \num{۲۶۵} سطر
از \num{۲۲۶۰} سطر آموزش می‌بیند.

\textbf{دوم، آزمونی که خود مداخله را بی‌اعتبار می‌کرد پیش از عدد ثبت شده بود و در جهت مثبت
عمل کرد.} ضمانت سازگاری رتبه‌ای \lr{CORAL} تنها وقتی برای هر نمونه برقرار است که سوگیری‌های
آموخته‌شده مرتب باشند، و هیچ چیز در پارامترسازی آن‌ها را مرتب نمی‌کند. هر ده تانسور دقیقاً
مرتب درآمدند. پس این نتیجهٔ منفی دربارهٔ \lr{CORAL} است، نه دربارهٔ پیاده‌سازی شکسته‌ای که نام
\lr{CORAL} را یدک می‌کشد.

\subsection{\lr{ConvNeXt}: نتیجهٔ منفی که کران دارد، نه نتیجهٔ منفی تمیز}

سطر \lr{ConvNeXt} دو صلاحیت دارد که در جدول جا نمی‌شوند و حذفشان گزارش را قوی‌تر از شواهد
نشان می‌دهد.

نخست، این اجرا روی دو تا از پنج تا انجام شد، روی \num{۹۱۰} تصویر برون‌تایی، و به همان صورت گزارش
می‌شود. دوم، و مهم‌تر: \lr{ConvNeXt} در یکی از آن دو تا هنگام پایان زمان‌بندی \emph{هنوز در
حال بهتر شدن بود}، در حالی که تای دیگر شانزده دوره پیش از پایان به فلات رسیده بود. تقسیم همان
مقایسه به تفکیک تا نشان می‌دهد کسری معنادار در تای بریده‌شده متمرکز است و تای همگرا
تفکیک‌ناپذیر است. این تقسیم پس‌نگر است و تعداد نمونه را نصف می‌کند، پس صلاحیت است، نه نتیجه.

\lr{ConvNeXt} با این کار تبرئه نمی‌شود — برآورد نقطه‌ای تای همگرا هم منفی است — ولی
\textbf{بخشی از منفی بودن این نتیجه به کوتاه بودن زمان‌بندی تای یک وابسته است، و این را صریح
می‌نویسیم}. معناداری مجموع تا حدی گزارشی دربارهٔ زمان‌بندی چهل‌دوره‌ای است، نه صرفاً دربارهٔ
معماری.

\textbf{سه چیز را باید کنار هم خواند، و هر سه در متن آمده‌اند تا خواننده مجبور نباشد حدس بزند.}
نخست، این نتیجه روی \textbf{دو تا از پنج تا} به دست آمده است، روی \num{۹۱۰} تصویر برون‌تایی، نه روی
\num{۲۲۶۰} تصویر کامل استخر. دوم، از آن دو تا، \textbf{یکی با زمان‌بندی بریده شد}: بهترین دورهٔ آن
دورهٔ سی‌ونهم از چهل بود و هنوز بالا می‌رفت. سوم، تفکیک همان مقایسه به تفکیک تا، کسری معنادار را
در همان تای بریده‌شده می‌گذارد و تای همگرا را تفکیک‌ناپذیر نشان می‌دهد.

\textbf{و تصمیم گرفتیم آن محاسبه را خرج نکنیم.} رفع این ابهام نیازمند اجرای دوبارهٔ
\lr{ConvNeXt} با زمان‌بندی بلندتر روی همان تاهاست: هر دوره \num{۲۱۲} ثانیه در برابر حدود \num{۴۲} ثانیه برای
\lr{DenseNet121}، یعنی حدود سیزده ساعت و نیم پردازنده گرافیکی برای دو تا. این هزینه \emph{آگاهانه}
پرداخت نشد، و دلیلش این است که هیچ پیامدی در این خط به عدد سرتیتر نمی‌رسد: مدل انتخاب‌شده در
\num{۰٫۸۹۵۴} ایستاده و کنترل این آزمایش در \num{۰٫۸۶۲۸} است، پس حتی نتیجهٔ کاملاً مطلوب هم \lr{ConvNeXt} را
به مدل برگزیده نمی‌رساند. به همین دلیل \textbf{تاهای باقی‌مانده هم اجرا نشدند}.

این را به‌عنوان یک محدودیت می‌نویسیم، نه به‌عنوان یک حاشیه. \textbf{یک نتیجهٔ منفیِ صریحاً
مشروط از یک نتیجهٔ منفیِ صاف‌شده قوی‌تر است}: خواننده دقیقاً می‌داند چه چیزی اندازه گرفته شد،
روی چه اندازه‌ای از داده، زیر کدام زمان‌بندی، و کدام محاسبه پرداخت نشد و چرا. ادعای ما این نیست
که \lr{ConvNeXt} معماری بدتری است؛ ادعای ما این است که \textbf{در بودجهٔ این پژوهش و در برابر
مدل انتخاب‌شده، این خط ارزش ادامه دادن نداشت} — و همین جمله، نه بیشتر، چیزی است که شواهد
پشتیبانی می‌کنند.

\section{واسنجی همتاشده: سابقهٔ عمل در دو جهت}
\label{sec:matched-record}

قاعدهٔ بخش \ref{sec:matched} می‌گوید هیچ اختلافی به بازنمایی نسبت داده نمی‌شود مگر آن‌که از
واسنجی همتاشده جان به در ببرد. قاعده‌ای که همیشه در یک جهت عمل کند آزمون نیست، پس سابقه‌اش
اینجا می‌آید.

\begin{table}[H]
\caption{هر باری که واسنجی همتاشده روی یک ادعا اعمال شد، از جمله جایی که چیزی را عوض نکرد.}
\centering
\begin{tabular}{|l|l|l|}
\hline
\textbf{ادعا} & \textbf{تفسیر نخست} & \textbf{پس از واسنجی همتاشده} \\ \hline
فروپاشی بازخوانی درجهٔ خفیف & محدودیت ظرفیت & \textbf{واسنجی} — تنها با جابه‌جایی نقاط برش \\ \hline
برتری \lr{EfficientNet} & بازنمایی بهتر & برتری باقی ماند، با اندازهٔ کوچک‌تر \\ \hline
سود \lr{LP-FT} & انتقال بهتر & از میان رفت؛ تفکیک‌ناپذیر \\ \hline
زیان \lr{ConvNeXt} روی ادم ماکولا & معماری بدتر & \textbf{علامت وارونه شد}؛ تفکیک‌ناپذیر \\ \hline
\end{tabular}
\label{tab:matched_record}
\end{table}

سطر آخر تازه‌ترین است و همان چیزی است که قاعده برای آن ساخته شد. زیر نقاط برشِ عرضه‌شده،
\lr{ConvNeXt} روی سر ادم ماکولا بدتر به نظر می‌رسید؛ زیر نقاط برش همتاشده بهتر. هر دو بازه
صفر را در بر می‌گیرند، پس خواندنِ درست «تفکیک‌ناپذیر» است — همان چیزی که سطر همتاشده گزارش
می‌کند. \textbf{اختلافی که با جابه‌جایی آستانه علامت عوض می‌کند، گزاره‌ای دربارهٔ آستانه
است، نه دربارهٔ بازنمایی.}

سطر نخست جهت مخالف را نشان می‌دهد و مهم‌ترین سطر جدول است: همان قاعده‌ای که سه ادعای مثبت را
از میان برد، یک نتیجهٔ به‌ظاهر منفی — کوری نسبت به درجهٔ خفیف — را به مسئله‌ای حل‌شدنی و
بی‌هزینه تبدیل کرد.

\section{بازنمایی در برابر ریزتنظیم}
\label{sec:retfound}

یک مدل بنیادین شبکیه با پیش‌آموزش خودنظارتی روی \num{۱٫۶} میلیون تصویر آزموده شد، در دو حالت.

با رمزگذار منجمد، بازنمایی آن \emph{بهتر} است: \lr{QWK} رتینوپاتی \num{۰٫۷۰۰۸} در برابر \num{۰٫۶۶۱۱} برای
\lr{ImageNet}، اختلاف \num{۰٫۰۴۰۱} با بازهٔ [\num{۰٫۰۱۲۲}، \num{۰٫۰۶۹۲}] که صفر را در بر نمی‌گیرد. روی سر ادم
ماکولا اختلاف \num{۰٫۰۴۳۴} با بازهٔ [-\num{۰٫۰۱۷۳}، \num{۰٫۱۰۱۸}] تفکیک‌ناپذیر است.

با ریزتنظیم کامل، همان مدل به خط لولهٔ خودمان \emph{می‌بازد}.

\textbf{دو سازوکار نامزد وجود دارد و این آزمایش نمی‌تواند آن‌ها را از هم جدا کند.} یکی این‌که
دادهٔ برچسب‌خوردهٔ درون‌دامنه در این مقیاس بر دادهٔ بی‌برچسب می‌چربد. دیگری این‌که ظرفیت مدل
در برابر اندازهٔ نمونه بیش از حد است. این دو به‌طور ساختاری درهم‌تنیده‌اند، چون مدل بنیادین
هم‌زمان ستون فقرات و پیکرهٔ پیش‌آموزش را عوض می‌کند. هر دو به‌عنوان فرضیه ثبت می‌شوند، نه
به‌عنوان نتیجه.

\section{آنچه این فصل ادعا نمی‌کند}
\label{sec:limits}

\textbf{هیچ‌یک از این عددها با کار منتشرشده قابل مقایسه نیست.} تقسیم این پژوهش اختصاصی و منجمد
است و بر پایهٔ خوشه‌بندی ادراکی ساخته شده تا نشت میان‌پیکره‌ای را ببندد. همین کار آن را با هر
تقسیم منتشرشده‌ای ناهم‌سنجه می‌کند. عددی که کنار عدد یک مقاله گذاشته شود، دو چیز متفاوت را
مقایسه می‌کند.

\textbf{عدد بیرونی از یک پیکره می‌آید، نه از چند پیکره.} \lr{APTOS} یک جابه‌جایی توزیع واقعی
است، اما یک جابه‌جایی است.

\textbf{سر ادم ماکولا روی \num{۵۱۶} تصویر از یک پیکره سنجیده می‌شود}، با \num{۵۱} نمونه در درجهٔ میانی.
بازهٔ اطمینان روی این سر پهن است و هر ادعای کلاسی روی آن باید با همین پهنا خوانده شود.

\textbf{و کارایی، سهم اصلی این پایان‌نامه نیست.} هشت ابطال از نه مداخله، و آنچه در فصل
\ref{sec:parta} می‌آید، نشان می‌دهند که آنچه این پژوهش دارد بیش از آن‌که یک عدد بهتر باشد، یک
رویهٔ قابل دفاع برای رسیدن به عدد است.

% گزارش Part A — آزمون اشتقاق هندسی درجهٔ ادم ماکولا
% همهٔ عددها از docs/generated/idrid_derivation_gate.json می‌آیند و با اسکریپت بازتولید می‌شوند.
% معیار پذیرش پیش از محاسبهٔ هر عدد در docs/PARTA_preregistration.md ثبت و کامیت شد.

\section{آزمون اشتقاق هندسی درجهٔ ادم ماکولا}
\label{sec:parta}

\subsection{پرسش و انگیزه}

درجهٔ ادم ماکولای دیابتی در مجموعهٔ \lr{IDRiD} تعریفی هندسی و قطعی دارد: درجهٔ صفر یعنی نبودِ
ترشح سخت؛ درجهٔ دو یعنی کوتاه‌ترین فاصلهٔ هر پیکسل ترشح سخت تا مرکز فووه‌آ کمتر یا مساوی یک قطر
دیسک نوری باشد؛ و هر حالت دیگر درجهٔ یک است.

اگر این تعریف واقعاً همان تابعی باشد که متخصص به کار می‌برد، آن‌گاه هر مجموعه‌ای که ماسک ترشح
سخت و مختصات فووه‌آ و مرز دیسک نوری داشته باشد، به یک منبع برچسب سه‌کلاسه تبدیل می‌شود. این
تنها مسیری است که در این پژوهش برای افزایش دادهٔ درجهٔ میانی یافتیم، و درجهٔ میانی همان کلاسی
است که کمبودش سقف کارایی سر ادم ماکولا را تعیین می‌کند (یافتهٔ \lr{F7}).

زیرمجموعهٔ قطعه‌بندی \lr{IDRiD} شرایط ایده‌آل این آزمون را فراهم می‌کند: ماسک پیکسلی ترشح سخت،
مختصات فووه‌آ از حاشیه‌نویسی متخصص، ماسک دیسک نوری، و مهم‌تر از همه، تفکیک ترشح سخت از ترشح نرم.
این تفکیک را پیش از آزمون بررسی کردیم و تأیید شد: ماسک‌ها در دو پوشهٔ جدا نگه‌داری می‌شوند،
۸۱ تصویر ماسک ترشح سخت دارند و ۴۰ تصویر ماسک ترشح نرم.

\subsection{معیار پذیرش، پیش از محاسبه}

معیار پیش از محاسبهٔ هر عدد نوشته و کامیت شد. تصمیم گرفتیم پنج چیز را بدون استثنا گزارش کنیم:
نرخ تطابق دقیق، ماتریس درهم‌ریختگی سه‌در‌سه، بازخوانی هر کلاس با \textbf{گزارش جداگانهٔ درجهٔ یک}،
تعداد تصویرهایی که آزمون واقعاً پوشش می‌دهد، و تکرار کل آزمون زیر سه تعریف مختلف از «یک قطر دیسک».

عمداً هیچ آستانهٔ عددی برای «قبولی» تعیین نکردیم. تعیین آستانه‌ای مانند ۸۰ درصد بحث را به
درستیِ خودِ آستانه منتقل می‌کند. به جای آن تعهد کردیم که سطر درجهٔ یک حکم را حمل کند، چون
اشتقاقی که نتواند درجهٔ میانی را جای‌گذاری کند برای هدفی که این کار را برانگیخته بی‌فایده است.

\subsection{دو محدودیت که آزمون را پیش از اجرا تضعیف کردند}

\subsubsection{مجموعهٔ \lr{IDRiD} هیچ نگاشتی میان شماره‌گذاری‌هایش منتشر نکرده است}

زیرمجموعهٔ قطعه‌بندی با \lr{IDRiD\_01} تا \lr{IDRiD\_81} شماره‌گذاری شده و زیرمجموعهٔ درجه‌بندی
با \lr{IDRiD\_001} تا \lr{IDRiD\_413} برای آموزش و \lr{IDRiD\_001} تا \lr{IDRiD\_103} برای آزمون.
این دو شماره‌گذاری به هم نگاشته نمی‌شوند و ناشر هیچ جدول تناظری منتشر نکرده است.

تناظر را از خودِ تصویرها بازسازی کردیم، با دو سیگنال مستقل که هر دو باید هم‌داستان می‌شدند:
همبستگی پیکسلی میان تصویر اصلی قطعه‌بندی و هر تصویر درجه‌بندی، با آستانهٔ بزرگ‌تر از ۰٫۹۹۹۹؛ و
فاصلهٔ مرکز ثقل ماسک دیسک نوری تا مختصات دیسک نوری منتشرشده در پروندهٔ مکان‌یابی، با آستانهٔ
کمتر از ۱۰۰ پیکسل.

هیچ‌کدام به تنهایی کافی نیست. تصویرهای ته‌چشم در مقیاس کلی به هم شبیه‌اند، پس دومین گزینهٔ
همبستگی در فاصلهٔ کمتر از ۰٫۰۱ قرار می‌گیرد؛ و دیسک نوری در بیشتر چشم‌ها در ناحیهٔ مشابهی است،
پس نزدیک‌ترین دیسک یکتا نیست. هم‌زمانی این دو قوی است: یک جفت غلط باید هم تطابق پیکسلی داشته
باشد و هم دیسک نوری هم‌مکان.

از ۸۱ تصویر، \textbf{۵۴ تصویر} تأیید شد و ۲۷ تصویر تأییدنشده ماند.

\subsubsection{زیرمجموعهٔ قطعه‌بندی برای «تصویرهای ارزشمند برای قطعه‌بندی» انتخاب شده است}

توزیع درجهٔ متخصص در ۵۴ تصویر قابل استفاده چنین است: \textbf{یک تصویر درجهٔ صفر، سه تصویر
درجهٔ یک، و ۵۰ تصویر درجهٔ دو}. علت روشن است: سازندگان مجموعه تصویرهایی را برای حاشیه‌نویسی
پیکسلی انتخاب کردند که ضایعه داشته باشند، و آن تصویرها تقریباً همه درجهٔ دو هستند.

پیامد این توزیع را باید صریح نوشت: \textbf{یک پیش‌بین ثابت که همیشه درجهٔ دو بگوید، ۹۲٫۶ درصد
تطابق می‌گیرد.} پس هر عدد تطابق کلی را باید در برابر این کف خواند، و آزمون در عمل تنها روی
\textbf{چهار تصویر} قدرت تفکیک دارد. خواننده نباید تصور کند اشتقاق روی ۵۴ تصویر آزموده شده است.

\subsection{نتیجه}

\begin{table}[h]
\centering
\caption{اشتقاق هندسی در برابر درجهٔ متخصص، زیر سه تعریف از یک قطر دیسک. تعداد نمونه‌ها در
پرانتز آمده است. سطر درجهٔ یک حکم را حمل می‌کند.}
\label{tab:parta}
\begin{tabular}{lcccc}
\hline
تعریف یک قطر دیسک & تطابق دقیق & درجهٔ ۰ (۱) & \textbf{درجهٔ ۱ (۳)} & درجهٔ ۲ (۵۰) \\
\hline
محور بزرگ           & ۹۶٫۳\٪ & ۰٫۰\٪ & \textbf{۶۶٫۷\٪} & ۱۰۰٫۰\٪ \\
قطر معادل مساحت     & ۹۶٫۳\٪ & ۰٫۰\٪ & \textbf{۶۶٫۷\٪} & ۱۰۰٫۰\٪ \\
محور کوچک           & ۹۶٫۳\٪ & ۰٫۰\٪ & \textbf{۶۶٫۷\٪} & ۱۰۰٫۰\٪ \\
\hline
پیش‌بین ثابت «درجهٔ ۲» & ۹۲٫۶\٪ & ۰٫۰\٪ & ۰٫۰\٪ & ۱۰۰٫۰\٪ \\
\hline
\end{tabular}
\end{table}

سه تعریف از قطر دیسک نتیجهٔ یکسان می‌دهند. این را پیش از آزمون به‌عنوان یک ابهام روش‌شناختی
مطرح کرده بودیم، چون دیسک نوری بیضی است و آستانه بر حسب «قطر» تعریف شده است. اکنون بسته است:
فاصلهٔ ترشح تا فووه‌آ در همهٔ نمونه‌ها از آستانه دور است — میان ۵۳ تا ۷۵۶ پیکسل در برابر قطرهای
۴۵۷ تا ۶۵۷ پیکسل — پس انتخاب تعریف حکم را عوض نمی‌کند.

\subsection{چهار تصویری که آزمون روی آن‌ها تفکیک دارد}

\begin{table}[h]
\centering
\caption{تنها تصویرهایی که درجهٔ متخصصشان دو نیست. ستون «نزدیک‌ترین لکه» اندازهٔ همان مؤلفهٔ
همبندی است که درجهٔ مشتق را تعیین می‌کند.}
\label{tab:parta_four}
\begin{tabular}{lccccc}
\hline
تصویر & ترشح سخت & نزدیک‌ترین لکه & فاصله تا فووه‌آ & مشتق & متخصص \\
\hline
\lr{IDRiD\_29} & ۲۳۴۳ پیکسل، ۵ لکه   & ۱۳۹ پیکسل & ۷۵۶ پیکسل & ۱ & ۱ \checkmark \\
\lr{IDRiD\_62} & ۴۴۴۲۷ پیکسل، ۶۸ لکه & ۶۵۷ پیکسل & ۵۶۰ پیکسل & ۱ & ۱ \checkmark \\
\lr{IDRiD\_40} & ۷۷۷۹ پیکسل، ۱۹ لکه  & \textbf{۷۴ پیکسل} & ۵۳ پیکسل & ۲ & ۱ \ding{55} \\
\lr{IDRiD\_52} & ۲۹۷۴ پیکسل، ۱۳ لکه  & \textbf{۳۴۹ پیکسل} & ۸۷ پیکسل & ۲ & ۰ \ding{55} \\
\hline
\end{tabular}
\end{table}

اشتقاق دو تصویر را درست و دو تصویر را غلط تعیین می‌کند، و هر دو خطا در جهتی است که اهمیت دارد:
تعریف می‌گوید درجهٔ دو، متخصص می‌گوید صفر یا یک.

\subsection{دو توضیح رقیب که پیش از نتیجه‌گیری رد شدند}

نتیجه‌گیری از دو تصویر ناسازگار نیازمند احتیاط است. دو توضیح رقیب وجود داشت که هر دو همین
الگو را می‌سازند، و هر دو را بررسی کردیم.

\textbf{توضیح نخست: تناظر غلط.} یک جفت اشتباه دقیقاً همین تصویر را می‌سازد — ماسکی پر از ترشح
کنار درجه‌ای بی‌ربط. برای رد آن، ماسک ترشح سخت و مختصات فووه‌آ را روی \emph{تصویر درجه‌بندی
متناظر} هم‌نهاد کردیم و پیکسل‌ها را مقایسه کردیم. میانگین قدرمطلق اختلاف برای هر چهار جفت میان
۰٫۶۵ تا ۰٫۷۵ از ۲۵۵ است، یعنی تنها نوفهٔ فشرده‌سازی \lr{JPEG}. بازرسی چشمی نیز نشان می‌دهد نشانگر
فووه‌آ روی ماکولا و دیسک نوری در سمت مورد انتظار قرار دارد. تناظر درست است.

\textbf{توضیح دوم: تعریف متفاوت.} شاید \lr{IDRiD} درجه را طور دیگری تعریف کرده باشد. متن منتشرشدهٔ
مجموعه را بررسی کردیم: تعریف دقیقاً همان است که پیاده کردیم.

\subsection{تشخیص نهایی}

پس تعریف همان است و تناظر درست است. خطا جای دیگری است، و دقیق‌تر و مفیدتر از «متخصص تعریف
دیگری به کار برده» است:

\textbf{قاعدهٔ کمینهٔ فاصله روی یک ماسک پیکسلی را کوچک‌ترین لکهٔ حاشیه‌نویسی‌شده تعیین می‌کند.}
در \lr{IDRiD\_40} لکه‌ای به اندازهٔ ۷۴ پیکسل در تصویری با ابعاد ۲۸۴۸ در ۴۲۸۸ درجه را از یک به دو
می‌برد. در هر دو خطا، لکهٔ تعیین‌کننده کوچک است و درست داخل آستانه قرار دارد؛ در هر دو مورد
درست، ترشح‌ها به‌روشنی بیرون از یک قطر دیسک‌اند.

واژه‌ای که بار معنایی را حمل می‌کند «آشکار» است. درجهٔ صفر بر پایهٔ نبودِ ترشح \emph{آشکار}
تعریف شده، یعنی آنچه بالینگر می‌بیند؛ اما ماسک قطعه‌بندی هر لکه‌ای را که یک حاشیه‌نویس در سطح
پیکسل بتواند بیابد علامت می‌زند. این‌ها دو شیء متفاوت‌اند، حاصل دو کوشش حاشیه‌نویسی جدا با دو
هدف جدا، و قاعدهٔ کمینهٔ فاصله اختلافشان را بیشینه می‌کند چون به کوچک‌ترینشان حساس است.

این یک نتیجهٔ روش‌شناختی عام است، نه ویژگی \lr{IDRiD}: \textbf{اشتقاق یک برچسب ردهای بالینی از
ماسک قطعه‌بندی با آستانهٔ فاصله، کفِ حاشیه‌نویسی ماسک را به مرز تصمیم خود تبدیل می‌کند.}

\subsection{پیامدها}

مسیر اشتقاق بسته شد. مجموعهٔ \lr{SUSTech-SYSU} ارزش مهندسی ندارد، چون دو ضعف شناخته‌شده‌اش —
ترشح‌ها به شکل جعبهٔ محصورکننده نه ماسک پیکسلی، و احتمال ادغام ترشح سخت و نرم — روی روشی سوار
می‌شود که در شرایط ایده‌آل هم شکست خورده است.

آیا آستانهٔ اندازه می‌تواند آن را ترمیم کند؟ تشخیص بالا چنین چیزی را پیشنهاد می‌کند، ولی این
داده نمی‌تواند به آن پاسخ دهد: چهار تصویر تفکیک‌کننده و دو خطا در اختیار است، پس هر آستانه‌ای
روی دو نمونه برازش می‌شود و روی هیچ نمونه‌ای اعتبارسنجی نمی‌شود. آن را به‌عنوان پرسش بعدی برای
کسی که مجموعه‌ای با توان آماری کافی دارد ثبت می‌کنیم، نه به‌عنوان چیزی که اینجا حل شود.

سرانجام، این یافته به یافتهٔ \lr{F2} سازوکار می‌دهد. پیش‌تر ثبت کرده بودیم که هیچ مجموعهٔ
سه‌کلاسهٔ دوم برای ادم ماکولا وجود ندارد. اکنون دلیلی برایش داریم: درجه از روی حاشیه‌نویسی‌هایی
که این حوزه منتشر می‌کند محاسبه‌پذیر نیست، پس ساختن چنین مجموعه‌ای به درجه‌بندی تازهٔ متخصص نیاز
دارد، نه به محاسبه روی برچسب‌های موجود.

% فصل ۵ — بحث، محدودیت‌ها و کارهای آینده
% همهٔ عددها از docs/generated/thesis_numbers.json می‌آیند.

\chapter{بحث}
\label{ch:discussion}

\section{آنچه این پژوهش نشان داد}

سه نتیجه، به ترتیب اهمیت.

\textbf{یکم، قاعدهٔ تصمیم بخش بزرگی از آن چیزی را توضیح می‌دهد که معمولاً به معماری نسبت داده
می‌شود.} این نتیجه از یک آزمایش نمی‌آید، از سابقهٔ عمل یک قاعده می‌آید: واسنجی همتاشده تا امروز
چهار ادعا را در \emph{دو جهت} برگردانده و دو بار علامت اختلاف را وارونه کرده است. تازه‌ترین
مورد در فصل \ref{ch:results} آمد — روی سر ادم ماکولا، زیر نقاط برشِ عرضه‌شده یک معماری بدتر
به نظر می‌رسید و زیر نقاط برش همتاشده بهتر، در حالی که هر دو بازه صفر را در بر می‌گیرند.

\textbf{اهمیت این نتیجه در جهت مثبتش است، نه منفی‌اش.} همان قاعده‌ای که ادعاهای مثبت را از میان
برد، کوری مدل نسبت به درجهٔ خفیف را از یک محدودیت ظرفیت به یک مسئلهٔ واسنجی تبدیل کرد: بدون
آموزش دوباره و بدون تغییر مدل، تنها با جابه‌جایی نقاط برش. \textbf{یک قاعده که فقط ادعاها را
رد کند، بدبینی است؛ قاعده‌ای که هم رد کند و هم بسازد، آزمون است.}

\textbf{دوم، سقف سر ادم ماکولا از داده می‌آید.} شش مداخلهٔ مستقل بازهٔ اطمینان را از صفر دور
نکردند، از جمله مداخله‌ای که مستقیماً ناحیهٔ تعریف‌کنندهٔ درجهٔ بالینی را جدا می‌کند. تنها
\textbf{\num{۵۱}} تصویر در کل این پژوهش درجهٔ میانی ادم ماکولا دارند.

این نتیجه به تنهایی می‌توانست دربارهٔ انتخاب پیکرهٔ ما باشد. \textbf{دو شاهد بیرونی آن راه را
می‌بندند.} یک پیکرهٔ ادم ماکولا که برای همین کار ساخته شده و \textbf{\num{۲۹۳۸}} تصویر دارد، باز هم
تنها \textbf{\num{۹٫۵}} درصد نمونه در درجهٔ میانی دارد \cite{tang2026mmrdr}: نسبت میانی ویژگی این
بیماری است، نه ویژگی مجموعهٔ ما. و یک سامانهٔ تجاری مستقر، در یک اجرا و روی همان تصویرها، برای
رتینوپاتی نیازمند ارجاع به توافق \textbf{\num{۰٫۸۱}} می‌رسد و برای ادم ماکولا به حساسیت
\textbf{\num{۲۶٫۵}} درصد و توافق \textbf{\num{۰٫۳۸}} \cite{duggal2025realworld}. \textbf{همان تصویرها،
همان اجرا، یک وظیفه کار می‌کند و دیگری نه.}

\textbf{سوم، تبار داده در این حوزه — و در این پایان‌نامه — مکانیزم لازم دارد، نه دقت.} این
نتیجه در آغاز کار پیش‌بینی نشده بود و در فصل \ref{ch:method} با پنج نمونه آمد که یکی از آن‌ها
خودِ ماست.

\section{سه تفسیری که شواهد ما اجازه نمی‌دهند}

\subsection{مدل بنیادین: دو سازوکار که از هم جدا نمی‌شوند}

با رمزگذار منجمد، بازنمایی مدل بنیادین شبکیه بهتر از \lr{ImageNet} است؛ با ریزتنظیم کامل، به
خط لولهٔ سادهٔ ما می‌بازد. دو توضیح وجود دارد. یکی این‌که دادهٔ برچسب‌خوردهٔ درون‌دامنه در این
مقیاس بر دادهٔ بی‌برچسب می‌چربد. دیگری این‌که ظرفیت مدل در برابر اندازهٔ نمونه بیش از حد است.

\textbf{این دو به‌طور ساختاری درهم‌تنیده‌اند و این آزمایش نمی‌تواند جدایشان کند}، چون مدل
بنیادین هم‌زمان ستون فقرات و پیکرهٔ پیش‌آموزش را عوض می‌کند. جدا کردنشان نیازمند اجرایی است که
تنها یکی را عوض کند، و آن اجرا انجام نشده. \textbf{هر دو به‌عنوان فرضیه ثبت می‌شوند، نه
به‌عنوان نتیجه.}

\subsection{\lr{ConvNeXt}: نتیجه‌ای که به زمان‌بندی مشروط است}

نتیجهٔ منفی \lr{ConvNeXt} روی دو تا از پنج تا به دست آمد و در یکی از آن دو تا، مدل هنگام پایان
زمان‌بندی هنوز در حال بهتر شدن بود. تفکیک به‌تفکیک تا، کسری معنادار را در همان تای بریده‌شده
می‌گذارد. \textbf{ادعای ما این نیست که این معماری بدتر است؛ ادعای ما این است که در بودجهٔ این
پژوهش و در برابر مدل انتخاب‌شده، این خط ارزش ادامه دادن نداشت.}

\subsection{ناهنجاری‌ای که تشخیصش داده نشد}

هنگام راستی‌آزمایی، در یکی از مقاله‌های مرجع دو عدد \lr{AUROC} — یکی درونی و یکی بیرونی، از دو
رژیم آموزشی متفاوت — تا سه رقم اعشار و \emph{شامل بازهٔ اطمینان} یکسان گزارش شده‌اند.
\textbf{مشاهده تأیید شده است؛ تشخیص نه.} از روی مطالب منتشرشده نتوانستیم تعیین کنیم که این
هم‌زمانی تصادف است یا خطای رونویسی، و \textbf{ادعای خطا در یک مقاله بر پایهٔ این شواهد را
مجاز نمی‌دانیم.} حل آن نیازمند جدول‌های تکمیلی است که در دسترس نبودند. ثبت می‌شود چون
\emph{نگفتنش} هم انتخابی است.

\section{محدودیت‌ها}

\subsection{ناهم‌سنجگی با کار منتشرشده}

\textbf{این جدی‌ترین محدودیت این پایان‌نامه است.} تقسیم این پژوهش اختصاصی و منجمد است و بر
پایهٔ خوشه‌بندی ادراکی ساخته شده تا نشت میان‌پیکره‌ای را ببندد. همین انتخاب — که برای درستی
لازم بود — هر عدد این متن را با هر عدد منتشرشده‌ای ناهم‌سنجه می‌کند.

\textbf{و این یک بده‌بستان است، نه یک نقص.} می‌شد از یک تقسیم استاندارد استفاده کرد و اعداد
قابل مقایسه گرفت؛ آنگاه اعداد آلوده به نشت می‌بودند. راهی که انتخاب شد، مقایسه‌پذیری را
فدای درستی می‌کند، و خواننده باید بداند کدام را در دست دارد.

\subsection{اعتبارسنجی بیرونی از یک پیکره می‌آید}

\lr{APTOS} یک جابه‌جایی توزیع واقعی است، اما \emph{یک} جابه‌جایی. برچسب‌هایش تک‌داورند و
نوفه‌دارتر، پس بخشی از هر تفاوت درون‌به‌بیرون نوفهٔ برچسب است، نه ناتوانی در تعمیم. ادعای
تعمیم‌پذیری عمومی از یک مجموعهٔ بیرونی برنمی‌آید.

\subsection{سر ادم ماکولا روی یک پیکره سنجیده می‌شود}

\textbf{\num{۵۱۶}} تصویر از یک منبع، با \textbf{\num{۵۱}} نمونه در درجهٔ میانی. بازهٔ اطمینان روی این سر
پهن است و هر ادعای کلاسی باید با همین پهنا خوانده شود. فصل \ref{sec:parta} نشان می‌دهد چرا
پیکرهٔ دومی افزوده نشد.

\subsection{بدون تصحیح مقایسه‌های چندگانه}

پنج کلاس ضرب در شرایط مقایسه‌شده یعنی به‌طور متوسط یک خانهٔ معنادار از سر تصادف انتظار می‌رود.
اعتبار در این متن از تکرار یک اختلاف در شرایط مختلف می‌آید، نه از یک مقدار احتمال.

\subsection{محدودیت‌های بیرون از دامنه}

سه مانع دیگر استقرار که فصل \ref{ch:intro} نام برد — هزینه‌اثربخشی، عدالت و سوگیری، و مسئولیت
حقوقی-پزشکی \cite{ruamviboonsuk2024review} — در این پژوهش بررسی نشده‌اند.
\textbf{هیچ ادعایی دربارهٔ آمادگی بالینی در این متن نیست.} این کار یک مطالعهٔ روش‌شناختی روی
داده‌های عمومی است، نه ارزیابی یک ابزار.

\subsection{و محدودیتی که از خودِ ما بود}

نسخه‌های نخست این متن ادعا می‌کردند عددهایشان تولیدشده است در حالی که با دست تایپ شده بودند، و
سه عدد غلط بود که یکی از آن‌ها هیچ مصنوعی نداشت و بازتولید نشد. سپس، هنگام نوشتنِ همین گزارش،
تشخیصی از روی مشخصات یک استاندارد به‌جای نگاه کردن به صفحهٔ چاپ‌شده گرفته شد و چیزی را «تصحیح»
کرد که غلط نبود. \textbf{هر دو در بخش \ref{sec:our-own-failure} آمده‌اند و از فهرست
محدودیت‌ها حذف نمی‌شوند چون تصحیح شده‌اند.} آنچه یافتنشان را ممکن کرد سازوکار بود نه دقت، و
همین نکته است.

\section{کارهای آینده}

\textbf{یکم، جدا کردن دو سازوکار مدل بنیادین.} یک اجرا که تنها ستون فقرات را عوض کند و پیکرهٔ
پیش‌آموزش را ثابت نگه دارد، یا برعکس. طراحی درستش در ادبیات وجود دارد
\cite{zhou2023retfound} و باید عیناً دنبال شود.

\textbf{دوم، منحنی توافق در برابر نقطهٔ عمل.} این پژوهش نشان داد آماره‌های توافق تجمیعی نسبت به
جابه‌جایی نقطهٔ عمل کم‌حس‌اند؛ آنچه ندارد، منحنی کامل این کم‌حسی روی دادهٔ ثابت است. این
تحلیل \textbf{به هیچ محاسبهٔ تازه‌ای نیاز ندارد} و از روی پیش‌بینی‌های آرشیوشده به دست می‌آید.

\textbf{سوم، آستانهٔ ارجاع به‌عنوان یک انتخاب، نه یک پیش‌فرض.} هدف حساسیت باید نخست انتخاب و
در برابر رویهٔ غربالگری توجیه شود، سپس روی استخر توسعه برازش شود و شکاف انتقال گزارش شود.

\textbf{چهارم، سر تک‌خروجی برای رتینوپاتی.} هرگز اندازه نگرفتیم که سر مشترک ادم ماکولا به سر
رتینوپاتی کمک می‌کند یا آسیب می‌زند. یک اجرا، یک تغییر.

\textbf{و پنجم، آنچه دیگر دنبال نمی‌شود.} خط معماری بسته است: از دو مقایسهٔ تمیز، اثر معماری
چند صدم واحد است و در هر دو جهت. مسیر ساختن دادهٔ سه‌کلاسهٔ بیشتر از روی ماسک‌های قطعه‌بندی هم
بسته است، و فصل \ref{sec:parta} نشان می‌دهد چرا. \textbf{هر منبع تازهٔ ادم ماکولا باید درجهٔ
سه‌کلاسهٔ بومی داشته باشد، نه درجهٔ قابل استخراج.}

\section{نتیجه‌گیری}

این پایان‌نامه با پرسش «کدام معماری بهتر است؟» آغاز نشد و با پاسخی به آن هم پایان نمی‌یابد. از
نه مداخلهٔ آزموده‌شده، هشت مداخله ابطال شدند. آنچه ماند یک عدد نیست، یک رویه است: نقاط برش پیش
از انتساب همتا می‌شوند، هر اجرا بیانیهٔ مصرف می‌نویسد، هر قاعده سابقهٔ عمل خود را در هر دو جهت
حمل می‌کند، و هر عدد متن به مصنوعی روی دیسک وصل است.

\textbf{و آزمون این رویه این بود که خطای خودمان را پیدا کرد.} یک فصل که تبار دادهٔ دیگران را
حسابرسی کند و شکست خودش را از همان جنس گزارش نکند، ادعایی می‌کند که خودش را زیر آن نمی‌برد.
این متن آن ادعا را نمی‌کند.


% ------------------------------------------------------------------ کتاب‌نامه
\cleardoublepage
\addcontentsline{toc}{chapter}{کتاب‌نامه}
% کتاب‌نامه لاتین است و باید در بستر چپ‌به‌راست چیده شود. بدون این، هر مدخل داخل
% بند راست‌به‌چپ می‌افتد و ترتیب واژه‌ها وارونه چاپ می‌شود — نام نویسندگان از آخر به
% اول. با چاپ آزمایشی بررسی شد، نه با فرض.
% عنوان کتاب‌نامه فارسی است و بدنه‌اش لاتین. محیط latin هر دو را لاتین می‌کند و عنوان
% فارسی در قلم لاتین ناپدید می‌شود، پس عنوان بیرون از محیط چاپ می‌شود و سرفصلِ خودکارِ
% محیط کتاب‌نامه برداشته می‌شود.
\chapter*{کتاب‌نامه}
\makeatletter
\renewenvironment{thebibliography}[1]
  {\list{\@biblabel{\@arabic\c@enumiv}}%
     {\settowidth\labelwidth{\@biblabel{#1}}%
      \leftmargin\labelwidth \advance\leftmargin\labelsep
      \@openbib@code \usecounter{enumiv}\let\p@enumiv\@empty
      \renewcommand\theenumiv{\@arabic\c@enumiv}}%
   \sloppy\clubpenalty4000\widowpenalty4000\sfcode`\.\@m}
  {\def\@noitemerr{\@latex@warning{Empty bibliography}}\endlist}
\makeatother
\begin{latin}
\bibliographystyle{unsrt}
\bibliography{refs}
\end{latin}

\end{document}
